\documentclass[10pt,a4paper]{amsart}
\usepackage[margin=1in]{geometry}
\usepackage{barnette-figs}
\begin{document}
\section*{1. two factors}
\begin{figure}[h!]\centering
\begin{tabular}{@{}c@{}}
\begin{tikzpicture}[x=.72cm,y=.72cm]
\path[use as bounding box] (-.6,-.7) rectangle (12.6,11);
\draw[base] (11.00000,5.00000) -- (8.98126,4.17880);
\draw[base] (11.00000,5.00000) -- (9.53553,1.46447);
\draw[base] (11.00000,5.00000) -- (9.53553,8.53553);
\draw[base] (8.98126,4.17880) -- (8.42742,2.55324);
\draw[base] (8.98126,4.17880) -- (7.51638,4.98317);
\draw[base] (9.53553,1.46447) -- (8.42742,2.55324);
\draw[base] (9.53553,1.46447) -- (6.00000,0.00000);
\draw[base] (9.53553,8.53553) -- (6.00000,10.00000);
\draw[base] (9.53553,8.53553) -- (7.52740,7.68875);
\draw[base] (8.42742,2.55324) -- (6.76545,2.01644);
\draw[base] (7.51638,4.98317) -- (7.06034,6.08414);
\draw[base] (7.51638,4.98317) -- (6.50752,4.68658);
\draw[base] (6.76545,2.01644) -- (6.00000,0.00000);
\draw[base] (6.76545,2.01644) -- (5.86894,3.49608);
\draw[base] (6.00000,0.00000) -- (2.46447,1.46447);
\draw[base] (6.00000,10.00000) -- (2.46447,8.53553);
\draw[base] (6.00000,10.00000) -- (5.98632,8.44657);
\draw[base] (7.52740,7.68875) -- (7.06034,6.08414);
\draw[base] (7.52740,7.68875) -- (5.98632,8.44657);
\draw[base] (2.46447,1.46447) -- (1.00000,5.00000);
\draw[base] (2.46447,1.46447) -- (3.16972,3.17873);
\draw[base] (1.00000,5.00000) -- (2.46447,8.53553);
\draw[base] (1.00000,5.00000) -- (2.71086,4.28652);
\draw[base] (3.16972,3.17873) -- (4.33384,3.78522);
\draw[base] (3.16972,3.17873) -- (2.71086,4.28652);
\draw[base] (2.46447,8.53553) -- (4.43156,7.65095);
\draw[base] (2.71086,4.28652) -- (3.96287,4.68083);
\draw[base] (4.43156,7.65095) -- (5.98632,8.44657);
\draw[base] (4.43156,7.65095) -- (4.84389,5.97076);
\draw[base] (7.06034,6.08414) -- (6.13725,5.58049);
\draw[base] (6.50752,4.68658) -- (5.86894,3.49608);
\draw[base] (6.50752,4.68658) -- (6.13725,5.58049);
\draw[base] (6.13725,5.58049) -- (4.84389,5.97076);
\draw[base] (5.86894,3.49608) -- (4.33384,3.78522);
\draw[base] (4.33384,3.78522) -- (3.96287,4.68083);
\draw[base] (4.84389,5.97076) -- (3.96287,4.68083);
\draw[ham] (11.00000,5.00000) -- (9.53553,1.46447);
\draw[ham] (11.00000,5.00000) -- (9.53553,8.53553);
\draw[ham] (8.98126,4.17880) -- (8.42742,2.55324);
\draw[ham] (8.98126,4.17880) -- (7.51638,4.98317);
\draw[ham] (8.42742,2.55324) -- (6.76545,2.01644);
\draw[ham] (9.53553,1.46447) -- (6.00000,0.00000);
\draw[ham] (9.53553,8.53553) -- (7.52740,7.68875);
\draw[ham] (6.00000,0.00000) -- (6.76545,2.01644);
\draw[ham] (2.46447,1.46447) -- (1.00000,5.00000);
\draw[ham] (2.46447,1.46447) -- (3.16972,3.17873);
\draw[ham] (1.00000,5.00000) -- (2.71086,4.28652);
\draw[ham] (2.46447,8.53553) -- (6.00000,10.00000);
\draw[ham] (2.46447,8.53553) -- (4.43156,7.65095);
\draw[ham] (6.00000,10.00000) -- (5.98632,8.44657);
\draw[ham] (7.51638,4.98317) -- (6.50752,4.68658);
\draw[ham] (7.52740,7.68875) -- (7.06034,6.08414);
\draw[ham] (7.06034,6.08414) -- (6.13725,5.58049);
\draw[ham] (6.50752,4.68658) -- (5.86894,3.49608);
\draw[ham] (5.86894,3.49608) -- (4.33384,3.78522);
\draw[ham] (5.98632,8.44657) -- (4.43156,7.65095);
\draw[ham] (6.13725,5.58049) -- (4.84389,5.97076);
\draw[ham] (3.16972,3.17873) -- (2.71086,4.28652);
\draw[ham] (4.33384,3.78522) -- (3.96287,4.68083);
\draw[ham] (4.84389,5.97076) -- (3.96287,4.68083);
\node[circle,draw,fill=black,inner sep=1.4pt] at (11.00000,5.00000) {};
\node[circle,draw,fill=white,inner sep=1.4pt] at (8.98126,4.17880) {};
\node[circle,draw,fill=white,inner sep=1.4pt] at (9.53553,1.46447) {};
\node[circle,draw,fill=white,inner sep=1.4pt] at (9.53553,8.53553) {};
\node[circle,draw,fill=black,inner sep=1.4pt] at (8.42742,2.55324) {};
\node[circle,draw,fill=black,inner sep=1.4pt] at (7.51638,4.98317) {};
\node[circle,draw,fill=white,inner sep=1.4pt] at (6.76545,2.01644) {};
\node[circle,draw,fill=black,inner sep=1.4pt] at (6.00000,0.00000) {};
\node[circle,draw,fill=black,inner sep=1.4pt] at (6.00000,10.00000) {};
\node[circle,draw,fill=black,inner sep=1.4pt] at (7.52740,7.68875) {};
\node[circle,draw,fill=white,inner sep=1.4pt] at (2.46447,1.46447) {};
\node[circle,draw,fill=black,inner sep=1.4pt] at (1.00000,5.00000) {};
\node[circle,draw,fill=black,inner sep=1.4pt] at (3.16972,3.17873) {};
\node[circle,draw,fill=white,inner sep=1.4pt] at (2.46447,8.53553) {};
\node[circle,draw,fill=white,inner sep=1.4pt] at (2.71086,4.28652) {};
\node[circle,draw,fill=black,inner sep=1.4pt] at (4.43156,7.65095) {};
\node[circle,draw,fill=white,inner sep=1.4pt] at (5.98632,8.44657) {};
\node[circle,draw,fill=white,inner sep=1.4pt] at (7.06034,6.08414) {};
\node[circle,draw,fill=white,inner sep=1.4pt] at (6.50752,4.68658) {};
\node[circle,draw,fill=black,inner sep=1.4pt] at (6.13725,5.58049) {};
\node[circle,draw,fill=black,inner sep=1.4pt] at (5.86894,3.49608) {};
\node[circle,draw,fill=white,inner sep=1.4pt] at (4.33384,3.78522) {};
\node[circle,draw,fill=white,inner sep=1.4pt] at (4.84389,5.97076) {};
\node[circle,draw,fill=black,inner sep=1.4pt] at (3.96287,4.68083) {};
\end{tikzpicture}\\[1mm]Two cycle components\\[4mm]\begin{tikzpicture}[x=.72cm,y=.72cm]
\path[use as bounding box] (-.6,-.7) rectangle (12.6,11);
\draw[base] (11.00000,5.00000) -- (8.98126,4.17880);
\draw[base] (11.00000,5.00000) -- (9.53553,1.46447);
\draw[base] (11.00000,5.00000) -- (9.53553,8.53553);
\draw[base] (8.98126,4.17880) -- (8.42742,2.55324);
\draw[base] (8.98126,4.17880) -- (7.51638,4.98317);
\draw[base] (9.53553,1.46447) -- (8.42742,2.55324);
\draw[base] (9.53553,1.46447) -- (6.00000,0.00000);
\draw[base] (9.53553,8.53553) -- (6.00000,10.00000);
\draw[base] (9.53553,8.53553) -- (7.52740,7.68875);
\draw[base] (8.42742,2.55324) -- (6.76545,2.01644);
\draw[base] (7.51638,4.98317) -- (7.06034,6.08414);
\draw[base] (7.51638,4.98317) -- (6.50752,4.68658);
\draw[base] (6.76545,2.01644) -- (6.00000,0.00000);
\draw[base] (6.76545,2.01644) -- (5.86894,3.49608);
\draw[base] (6.00000,0.00000) -- (2.46447,1.46447);
\draw[base] (6.00000,10.00000) -- (2.46447,8.53553);
\draw[base] (6.00000,10.00000) -- (5.98632,8.44657);
\draw[base] (7.52740,7.68875) -- (7.06034,6.08414);
\draw[base] (7.52740,7.68875) -- (5.98632,8.44657);
\draw[base] (2.46447,1.46447) -- (1.00000,5.00000);
\draw[base] (2.46447,1.46447) -- (3.16972,3.17873);
\draw[base] (1.00000,5.00000) -- (2.46447,8.53553);
\draw[base] (1.00000,5.00000) -- (2.71086,4.28652);
\draw[base] (3.16972,3.17873) -- (4.33384,3.78522);
\draw[base] (3.16972,3.17873) -- (2.71086,4.28652);
\draw[base] (2.46447,8.53553) -- (4.43156,7.65095);
\draw[base] (2.71086,4.28652) -- (3.96287,4.68083);
\draw[base] (4.43156,7.65095) -- (5.98632,8.44657);
\draw[base] (4.43156,7.65095) -- (4.84389,5.97076);
\draw[base] (7.06034,6.08414) -- (6.13725,5.58049);
\draw[base] (6.50752,4.68658) -- (5.86894,3.49608);
\draw[base] (6.50752,4.68658) -- (6.13725,5.58049);
\draw[base] (6.13725,5.58049) -- (4.84389,5.97076);
\draw[base] (5.86894,3.49608) -- (4.33384,3.78522);
\draw[base] (4.33384,3.78522) -- (3.96287,4.68083);
\draw[base] (4.84389,5.97076) -- (3.96287,4.68083);
\draw[ham] (11.00000,5.00000) -- (9.53553,1.46447);
\draw[ham] (11.00000,5.00000) -- (9.53553,8.53553);
\draw[ham] (8.98126,4.17880) -- (8.42742,2.55324);
\draw[ham] (8.98126,4.17880) -- (7.51638,4.98317);
\draw[ham] (8.42742,2.55324) -- (9.53553,1.46447);
\draw[ham] (9.53553,8.53553) -- (6.00000,10.00000);
\draw[ham] (6.00000,0.00000) -- (2.46447,1.46447);
\draw[ham] (6.00000,0.00000) -- (6.76545,2.01644);
\draw[ham] (2.46447,1.46447) -- (3.16972,3.17873);
\draw[ham] (1.00000,5.00000) -- (2.46447,8.53553);
\draw[ham] (1.00000,5.00000) -- (2.71086,4.28652);
\draw[ham] (2.46447,8.53553) -- (4.43156,7.65095);
\draw[ham] (6.00000,10.00000) -- (5.98632,8.44657);
\draw[ham] (7.51638,4.98317) -- (6.50752,4.68658);
\draw[ham] (7.52740,7.68875) -- (7.06034,6.08414);
\draw[ham] (7.52740,7.68875) -- (5.98632,8.44657);
\draw[ham] (7.06034,6.08414) -- (6.13725,5.58049);
\draw[ham] (6.76545,2.01644) -- (5.86894,3.49608);
\draw[ham] (6.50752,4.68658) -- (5.86894,3.49608);
\draw[ham] (6.13725,5.58049) -- (4.84389,5.97076);
\draw[ham] (3.16972,3.17873) -- (4.33384,3.78522);
\draw[ham] (4.33384,3.78522) -- (3.96287,4.68083);
\draw[ham] (4.43156,7.65095) -- (4.84389,5.97076);
\draw[ham] (2.71086,4.28652) -- (3.96287,4.68083);
\node[circle,draw,fill=black,inner sep=1.4pt] at (11.00000,5.00000) {};
\node[circle,draw,fill=white,inner sep=1.4pt] at (8.98126,4.17880) {};
\node[circle,draw,fill=white,inner sep=1.4pt] at (9.53553,1.46447) {};
\node[circle,draw,fill=white,inner sep=1.4pt] at (9.53553,8.53553) {};
\node[circle,draw,fill=black,inner sep=1.4pt] at (8.42742,2.55324) {};
\node[circle,draw,fill=black,inner sep=1.4pt] at (7.51638,4.98317) {};
\node[circle,draw,fill=white,inner sep=1.4pt] at (6.76545,2.01644) {};
\node[circle,draw,fill=black,inner sep=1.4pt] at (6.00000,0.00000) {};
\node[circle,draw,fill=black,inner sep=1.4pt] at (6.00000,10.00000) {};
\node[circle,draw,fill=black,inner sep=1.4pt] at (7.52740,7.68875) {};
\node[circle,draw,fill=white,inner sep=1.4pt] at (2.46447,1.46447) {};
\node[circle,draw,fill=black,inner sep=1.4pt] at (1.00000,5.00000) {};
\node[circle,draw,fill=black,inner sep=1.4pt] at (3.16972,3.17873) {};
\node[circle,draw,fill=white,inner sep=1.4pt] at (2.46447,8.53553) {};
\node[circle,draw,fill=white,inner sep=1.4pt] at (2.71086,4.28652) {};
\node[circle,draw,fill=black,inner sep=1.4pt] at (4.43156,7.65095) {};
\node[circle,draw,fill=white,inner sep=1.4pt] at (5.98632,8.44657) {};
\node[circle,draw,fill=white,inner sep=1.4pt] at (7.06034,6.08414) {};
\node[circle,draw,fill=white,inner sep=1.4pt] at (6.50752,4.68658) {};
\node[circle,draw,fill=black,inner sep=1.4pt] at (6.13725,5.58049) {};
\node[circle,draw,fill=black,inner sep=1.4pt] at (5.86894,3.49608) {};
\node[circle,draw,fill=white,inner sep=1.4pt] at (4.33384,3.78522) {};
\node[circle,draw,fill=white,inner sep=1.4pt] at (4.84389,5.97076) {};
\node[circle,draw,fill=black,inner sep=1.4pt] at (3.96287,4.68083) {};
\end{tikzpicture}\\[1mm]One spanning cycle
\end{tabular}
\caption{Both bold subgraphs have degree two at every vertex. Only the lower one is connected. These are the same 24-vertex graph, dual to the running triangulation.}
\end{figure}
\clearpage
\section*{2. dual}
\begin{figure}[h!]\centering
\begin{tikzpicture}[x=0.72cm,y=0.72cm]
\path[use as bounding box] (-.6,-.7) rectangle (12.6,11);
\fill[black!9] (0.00000,0.00000) -- (1.32113,1.43991) -- (4.04994,0.42777) -- cycle;
\fill[black!9] (6.00000,10.39200) -- (2.69330,2.96621) -- (1.32113,1.43991) -- cycle;
\fill[black!9] (6.00000,10.39200) -- (6.18899,7.50100) -- (4.51516,5.61611) -- cycle;
\fill[black!9] (6.00000,10.39200) -- (10.01414,2.60858) -- (8.02049,5.26608) -- cycle;
\fill[black!9] (12.00000,0.00000) -- (4.04994,0.42777) -- (9.22826,0.75499) -- cycle;
\fill[black!9] (12.00000,0.00000) -- (7.28809,3.25717) -- (10.01414,2.60858) -- cycle;
\fill[black!9] (4.04994,0.42777) -- (2.69330,2.96621) -- (6.24482,1.48372) -- cycle;
\fill[black!9] (9.22826,0.75499) -- (6.24482,1.48372) -- (7.28809,3.25717) -- cycle;
\fill[black!9] (2.69330,2.96621) -- (4.51516,5.61611) -- (5.23909,3.48455) -- cycle;
\fill[black!9] (8.02049,5.26608) -- (7.28809,3.25717) -- (6.18899,7.50100) -- cycle;
\fill[black!9] (6.18899,7.50100) -- (6.24482,1.48372) -- (5.23909,3.48455) -- cycle;
\draw[base] (0.00000,0.00000) -- (6.00000,10.39200);
\draw[base] (0.00000,0.00000) -- (12.00000,0.00000);
\draw[base] (0.00000,0.00000) -- (4.04994,0.42777);
\draw[base] (0.00000,0.00000) -- (1.32113,1.43991);
\draw[base] (6.00000,10.39200) -- (12.00000,0.00000);
\draw[base] (6.00000,10.39200) -- (1.32113,1.43991);
\draw[base] (6.00000,10.39200) -- (10.01414,2.60858);
\draw[base] (6.00000,10.39200) -- (2.69330,2.96621);
\draw[base] (6.00000,10.39200) -- (8.02049,5.26608);
\draw[base] (6.00000,10.39200) -- (4.51516,5.61611);
\draw[base] (6.00000,10.39200) -- (6.18899,7.50100);
\draw[base] (12.00000,0.00000) -- (4.04994,0.42777);
\draw[base] (12.00000,0.00000) -- (10.01414,2.60858);
\draw[base] (12.00000,0.00000) -- (9.22826,0.75499);
\draw[base] (12.00000,0.00000) -- (7.28809,3.25717);
\draw[base] (4.04994,0.42777) -- (1.32113,1.43991);
\draw[base] (4.04994,0.42777) -- (9.22826,0.75499);
\draw[base] (4.04994,0.42777) -- (2.69330,2.96621);
\draw[base] (4.04994,0.42777) -- (6.24482,1.48372);
\draw[base] (1.32113,1.43991) -- (2.69330,2.96621);
\draw[base] (10.01414,2.60858) -- (8.02049,5.26608);
\draw[base] (10.01414,2.60858) -- (7.28809,3.25717);
\draw[base] (9.22826,0.75499) -- (7.28809,3.25717);
\draw[base] (9.22826,0.75499) -- (6.24482,1.48372);
\draw[base] (2.69330,2.96621) -- (4.51516,5.61611);
\draw[base] (2.69330,2.96621) -- (5.23909,3.48455);
\draw[base] (2.69330,2.96621) -- (6.24482,1.48372);
\draw[base] (8.02049,5.26608) -- (7.28809,3.25717);
\draw[base] (8.02049,5.26608) -- (6.18899,7.50100);
\draw[base] (7.28809,3.25717) -- (6.18899,7.50100);
\draw[base] (7.28809,3.25717) -- (6.24482,1.48372);
\draw[base] (4.51516,5.61611) -- (6.18899,7.50100);
\draw[base] (4.51516,5.61611) -- (5.23909,3.48455);
\draw[base] (6.18899,7.50100) -- (5.23909,3.48455);
\draw[base] (6.18899,7.50100) -- (6.24482,1.48372);
\draw[base] (5.23909,3.48455) -- (6.24482,1.48372);
\fill (0.00000,0.00000) circle[radius=1.1pt];
\draw[line width=.7pt,fill=white] (0.00000,0.00000) circle[radius=2.4pt];
\node[fill=white,inner sep=.5pt,anchor=south] at (0.00000,0.10000) {$0$};
\fill (12.00000,0.00000) circle[radius=1.1pt];
\draw[line width=.7pt,fill=white] (12.00000,0.00000) circle[radius=2.4pt];
\node[fill=white,inner sep=.5pt,anchor=south] at (12.00000,0.10000) {$2$};
\fill (6.00000,10.39200) circle[radius=1.1pt];
\draw[line width=.7pt,fill=white] (6.00000,10.39200) circle[radius=2.4pt];
\node[fill=white,inner sep=.5pt,anchor=south] at (6.00000,10.49200) {$1$};
\fill (4.04994,0.42777) circle[radius=1.1pt];
\node[fill=white,inner sep=.5pt,anchor=south] at (4.04994,0.52777) {$3$};
\fill (1.32113,1.43991) circle[radius=1.1pt];
\node[fill=white,inner sep=.5pt,anchor=south] at (1.32113,1.53991) {$4$};
\fill (10.01414,2.60858) circle[radius=1.1pt];
\node[fill=white,inner sep=.5pt,anchor=south] at (10.01414,2.70858) {$5$};
\fill (9.22826,0.75499) circle[radius=1.1pt];
\node[fill=white,inner sep=.5pt,anchor=south] at (9.22826,0.85499) {$6$};
\fill (2.69330,2.96621) circle[radius=1.1pt];
\node[fill=white,inner sep=.5pt,anchor=south] at (2.69330,3.06621) {$7$};
\fill (8.02049,5.26608) circle[radius=1.1pt];
\node[fill=white,inner sep=.5pt,anchor=south] at (8.02049,5.36608) {$8$};
\fill (7.28809,3.25717) circle[radius=1.1pt];
\node[fill=white,inner sep=.5pt,anchor=south] at (7.28809,3.35717) {$9$};
\fill (4.51516,5.61611) circle[radius=1.1pt];
\node[fill=white,inner sep=.5pt,anchor=south] at (4.51516,5.71611) {$10$};
\fill (6.18899,7.50100) circle[radius=1.1pt];
\node[fill=white,inner sep=.5pt,anchor=south] at (6.18899,7.60100) {$11$};
\fill (5.23909,3.48455) circle[radius=1.1pt];
\node[fill=white,inner sep=.5pt,anchor=south] at (5.23909,3.58455) {$12$};
\fill (6.24482,1.48372) circle[radius=1.1pt];
\node[fill=white,inner sep=.5pt,anchor=south] at (6.24482,1.58372) {$13$};
\node[anchor=north] at (6,-.25) {$t_0$ is the outer black face};
\end{tikzpicture}
\caption{The running Eulerian triangulation has 14 vertices and 24 faces. Grey interiors denote black faces; the unbounded face is also black. Hollow vertices are the three roots. Vertex numbers remain fixed throughout.}
\end{figure}
\clearpage
\section*{3. cycle trees}
\begin{figure}[h!]\centering
\begin{tabular}{@{}c@{}}
\begin{tikzpicture}[x=0.72cm,y=0.72cm]
\path[use as bounding box] (-.6,-.7) rectangle (12.6,11);
\fill[black!9] (0.00000,0.00000) -- (1.32113,1.43991) -- (4.04994,0.42777) -- cycle;
\fill[black!9] (6.00000,10.39200) -- (2.69330,2.96621) -- (1.32113,1.43991) -- cycle;
\fill[black!9] (6.00000,10.39200) -- (6.18899,7.50100) -- (4.51516,5.61611) -- cycle;
\fill[black!9] (6.00000,10.39200) -- (10.01414,2.60858) -- (8.02049,5.26608) -- cycle;
\fill[black!9] (12.00000,0.00000) -- (4.04994,0.42777) -- (9.22826,0.75499) -- cycle;
\fill[black!9] (12.00000,0.00000) -- (7.28809,3.25717) -- (10.01414,2.60858) -- cycle;
\fill[black!9] (4.04994,0.42777) -- (2.69330,2.96621) -- (6.24482,1.48372) -- cycle;
\fill[black!9] (9.22826,0.75499) -- (6.24482,1.48372) -- (7.28809,3.25717) -- cycle;
\fill[black!9] (2.69330,2.96621) -- (4.51516,5.61611) -- (5.23909,3.48455) -- cycle;
\fill[black!9] (8.02049,5.26608) -- (7.28809,3.25717) -- (6.18899,7.50100) -- cycle;
\fill[black!9] (6.18899,7.50100) -- (6.24482,1.48372) -- (5.23909,3.48455) -- cycle;
\draw[base] (0.00000,0.00000) -- (6.00000,10.39200);
\draw[base] (0.00000,0.00000) -- (12.00000,0.00000);
\draw[base] (0.00000,0.00000) -- (4.04994,0.42777);
\draw[base] (0.00000,0.00000) -- (1.32113,1.43991);
\draw[base] (6.00000,10.39200) -- (12.00000,0.00000);
\draw[base] (6.00000,10.39200) -- (1.32113,1.43991);
\draw[base] (6.00000,10.39200) -- (10.01414,2.60858);
\draw[base] (6.00000,10.39200) -- (2.69330,2.96621);
\draw[base] (6.00000,10.39200) -- (8.02049,5.26608);
\draw[base] (6.00000,10.39200) -- (4.51516,5.61611);
\draw[base] (6.00000,10.39200) -- (6.18899,7.50100);
\draw[base] (12.00000,0.00000) -- (4.04994,0.42777);
\draw[base] (12.00000,0.00000) -- (10.01414,2.60858);
\draw[base] (12.00000,0.00000) -- (9.22826,0.75499);
\draw[base] (12.00000,0.00000) -- (7.28809,3.25717);
\draw[base] (4.04994,0.42777) -- (1.32113,1.43991);
\draw[base] (4.04994,0.42777) -- (9.22826,0.75499);
\draw[base] (4.04994,0.42777) -- (2.69330,2.96621);
\draw[base] (4.04994,0.42777) -- (6.24482,1.48372);
\draw[base] (1.32113,1.43991) -- (2.69330,2.96621);
\draw[base] (10.01414,2.60858) -- (8.02049,5.26608);
\draw[base] (10.01414,2.60858) -- (7.28809,3.25717);
\draw[base] (9.22826,0.75499) -- (7.28809,3.25717);
\draw[base] (9.22826,0.75499) -- (6.24482,1.48372);
\draw[base] (2.69330,2.96621) -- (4.51516,5.61611);
\draw[base] (2.69330,2.96621) -- (5.23909,3.48455);
\draw[base] (2.69330,2.96621) -- (6.24482,1.48372);
\draw[base] (8.02049,5.26608) -- (7.28809,3.25717);
\draw[base] (8.02049,5.26608) -- (6.18899,7.50100);
\draw[base] (7.28809,3.25717) -- (6.18899,7.50100);
\draw[base] (7.28809,3.25717) -- (6.24482,1.48372);
\draw[base] (4.51516,5.61611) -- (6.18899,7.50100);
\draw[base] (4.51516,5.61611) -- (5.23909,3.48455);
\draw[base] (6.18899,7.50100) -- (5.23909,3.48455);
\draw[base] (6.18899,7.50100) -- (6.24482,1.48372);
\draw[base] (5.23909,3.48455) -- (6.24482,1.48372);
\draw[pedge] (0.00000,0.00000) -- (12.00000,0.00000);
\draw[pedge] (4.04994,0.42777) -- (1.32113,1.43991);
\draw[pedge] (6.00000,10.39200) -- (1.32113,1.43991);
\draw[pedge] (6.00000,10.39200) -- (4.51516,5.61611);
\draw[pedge] (6.00000,10.39200) -- (8.02049,5.26608);
\draw[pedge] (12.00000,0.00000) -- (9.22826,0.75499);
\draw[pedge] (12.00000,0.00000) -- (10.01414,2.60858);
\draw[pedge] (2.69330,2.96621) -- (6.24482,1.48372);
\draw[pedge] (9.22826,0.75499) -- (6.24482,1.48372);
\draw[pedge] (4.51516,5.61611) -- (5.23909,3.48455);
\draw[pedge] (8.02049,5.26608) -- (7.28809,3.25717);
\draw[pedge] (6.18899,7.50100) -- (6.24482,1.48372);
\node[circle,fill=stateblue,inner sep=2.3pt] at (0.00000,0.00000) {};
\node[circle,fill=stateblue,inner sep=2.3pt] at (12.00000,0.00000) {};
\node[circle,fill=stateblue,inner sep=2.3pt] at (10.01414,2.60858) {};
\node[circle,fill=stateblue,inner sep=2.3pt] at (9.22826,0.75499) {};
\node[circle,fill=stateblue,inner sep=2.3pt] at (2.69330,2.96621) {};
\node[circle,fill=stateblue,inner sep=2.3pt] at (6.18899,7.50100) {};
\node[circle,fill=stateblue,inner sep=2.3pt] at (6.24482,1.48372) {};
\node[rectangle,fill=pairorange,inner sep=2.3pt] at (6.00000,10.39200) {};
\node[rectangle,fill=pairorange,inner sep=2.3pt] at (4.04994,0.42777) {};
\node[rectangle,fill=pairorange,inner sep=2.3pt] at (1.32113,1.43991) {};
\node[rectangle,fill=pairorange,inner sep=2.3pt] at (8.02049,5.26608) {};
\node[rectangle,fill=pairorange,inner sep=2.3pt] at (7.28809,3.25717) {};
\node[rectangle,fill=pairorange,inner sep=2.3pt] at (4.51516,5.61611) {};
\node[rectangle,fill=pairorange,inner sep=2.3pt] at (5.23909,3.48455) {};
\draw[line width=.7pt,fill=white] (0.00000,0.00000) circle[radius=2.4pt];
\draw[line width=.7pt,fill=white] (12.00000,0.00000) circle[radius=2.4pt];
\draw[line width=.7pt,fill=white] (6.00000,10.39200) circle[radius=2.4pt];
\node[anchor=north] at (6,-.25) {$t_0$ is the outer black face};
\end{tikzpicture}\\[1mm]Two induced trees in $T$\\[4mm]\begin{tikzpicture}[x=.72cm,y=.72cm]
\path[use as bounding box] (-.6,-.7) rectangle (12.6,11);
\draw[base] (11.00000,5.00000) -- (8.98126,4.17880);
\draw[base] (11.00000,5.00000) -- (9.53553,1.46447);
\draw[base] (11.00000,5.00000) -- (9.53553,8.53553);
\draw[base] (8.98126,4.17880) -- (8.42742,2.55324);
\draw[base] (8.98126,4.17880) -- (7.51638,4.98317);
\draw[base] (9.53553,1.46447) -- (8.42742,2.55324);
\draw[base] (9.53553,1.46447) -- (6.00000,0.00000);
\draw[base] (9.53553,8.53553) -- (6.00000,10.00000);
\draw[base] (9.53553,8.53553) -- (7.52740,7.68875);
\draw[base] (8.42742,2.55324) -- (6.76545,2.01644);
\draw[base] (7.51638,4.98317) -- (7.06034,6.08414);
\draw[base] (7.51638,4.98317) -- (6.50752,4.68658);
\draw[base] (6.76545,2.01644) -- (6.00000,0.00000);
\draw[base] (6.76545,2.01644) -- (5.86894,3.49608);
\draw[base] (6.00000,0.00000) -- (2.46447,1.46447);
\draw[base] (6.00000,10.00000) -- (2.46447,8.53553);
\draw[base] (6.00000,10.00000) -- (5.98632,8.44657);
\draw[base] (7.52740,7.68875) -- (7.06034,6.08414);
\draw[base] (7.52740,7.68875) -- (5.98632,8.44657);
\draw[base] (2.46447,1.46447) -- (1.00000,5.00000);
\draw[base] (2.46447,1.46447) -- (3.16972,3.17873);
\draw[base] (1.00000,5.00000) -- (2.46447,8.53553);
\draw[base] (1.00000,5.00000) -- (2.71086,4.28652);
\draw[base] (3.16972,3.17873) -- (4.33384,3.78522);
\draw[base] (3.16972,3.17873) -- (2.71086,4.28652);
\draw[base] (2.46447,8.53553) -- (4.43156,7.65095);
\draw[base] (2.71086,4.28652) -- (3.96287,4.68083);
\draw[base] (4.43156,7.65095) -- (5.98632,8.44657);
\draw[base] (4.43156,7.65095) -- (4.84389,5.97076);
\draw[base] (7.06034,6.08414) -- (6.13725,5.58049);
\draw[base] (6.50752,4.68658) -- (5.86894,3.49608);
\draw[base] (6.50752,4.68658) -- (6.13725,5.58049);
\draw[base] (6.13725,5.58049) -- (4.84389,5.97076);
\draw[base] (5.86894,3.49608) -- (4.33384,3.78522);
\draw[base] (4.33384,3.78522) -- (3.96287,4.68083);
\draw[base] (4.84389,5.97076) -- (3.96287,4.68083);
\draw[ham] (11.00000,5.00000) -- (9.53553,1.46447);
\draw[ham] (11.00000,5.00000) -- (9.53553,8.53553);
\draw[ham] (8.98126,4.17880) -- (8.42742,2.55324);
\draw[ham] (8.98126,4.17880) -- (7.51638,4.98317);
\draw[ham] (8.42742,2.55324) -- (9.53553,1.46447);
\draw[ham] (9.53553,8.53553) -- (6.00000,10.00000);
\draw[ham] (6.00000,0.00000) -- (2.46447,1.46447);
\draw[ham] (6.00000,0.00000) -- (6.76545,2.01644);
\draw[ham] (2.46447,1.46447) -- (3.16972,3.17873);
\draw[ham] (1.00000,5.00000) -- (2.46447,8.53553);
\draw[ham] (1.00000,5.00000) -- (2.71086,4.28652);
\draw[ham] (2.46447,8.53553) -- (4.43156,7.65095);
\draw[ham] (6.00000,10.00000) -- (5.98632,8.44657);
\draw[ham] (7.51638,4.98317) -- (6.50752,4.68658);
\draw[ham] (7.52740,7.68875) -- (7.06034,6.08414);
\draw[ham] (7.52740,7.68875) -- (5.98632,8.44657);
\draw[ham] (7.06034,6.08414) -- (6.13725,5.58049);
\draw[ham] (6.76545,2.01644) -- (5.86894,3.49608);
\draw[ham] (6.50752,4.68658) -- (5.86894,3.49608);
\draw[ham] (6.13725,5.58049) -- (4.84389,5.97076);
\draw[ham] (3.16972,3.17873) -- (4.33384,3.78522);
\draw[ham] (4.33384,3.78522) -- (3.96287,4.68083);
\draw[ham] (4.43156,7.65095) -- (4.84389,5.97076);
\draw[ham] (2.71086,4.28652) -- (3.96287,4.68083);
\node[circle,draw,fill=black,inner sep=1.4pt] at (11.00000,5.00000) {};
\node[circle,draw,fill=white,inner sep=1.4pt] at (8.98126,4.17880) {};
\node[circle,draw,fill=white,inner sep=1.4pt] at (9.53553,1.46447) {};
\node[circle,draw,fill=white,inner sep=1.4pt] at (9.53553,8.53553) {};
\node[circle,draw,fill=black,inner sep=1.4pt] at (8.42742,2.55324) {};
\node[circle,draw,fill=black,inner sep=1.4pt] at (7.51638,4.98317) {};
\node[circle,draw,fill=white,inner sep=1.4pt] at (6.76545,2.01644) {};
\node[circle,draw,fill=black,inner sep=1.4pt] at (6.00000,0.00000) {};
\node[circle,draw,fill=black,inner sep=1.4pt] at (6.00000,10.00000) {};
\node[circle,draw,fill=black,inner sep=1.4pt] at (7.52740,7.68875) {};
\node[circle,draw,fill=white,inner sep=1.4pt] at (2.46447,1.46447) {};
\node[circle,draw,fill=black,inner sep=1.4pt] at (1.00000,5.00000) {};
\node[circle,draw,fill=black,inner sep=1.4pt] at (3.16972,3.17873) {};
\node[circle,draw,fill=white,inner sep=1.4pt] at (2.46447,8.53553) {};
\node[circle,draw,fill=white,inner sep=1.4pt] at (2.71086,4.28652) {};
\node[circle,draw,fill=black,inner sep=1.4pt] at (4.43156,7.65095) {};
\node[circle,draw,fill=white,inner sep=1.4pt] at (5.98632,8.44657) {};
\node[circle,draw,fill=white,inner sep=1.4pt] at (7.06034,6.08414) {};
\node[circle,draw,fill=white,inner sep=1.4pt] at (6.50752,4.68658) {};
\node[circle,draw,fill=black,inner sep=1.4pt] at (6.13725,5.58049) {};
\node[circle,draw,fill=black,inner sep=1.4pt] at (5.86894,3.49608) {};
\node[circle,draw,fill=white,inner sep=1.4pt] at (4.33384,3.78522) {};
\node[circle,draw,fill=white,inner sep=1.4pt] at (4.84389,5.97076) {};
\node[circle,draw,fill=black,inner sep=1.4pt] at (3.96287,4.68083) {};
\end{tikzpicture}\\[1mm]Complementary dual cycle in $G$
\end{tabular}
\caption{Every edge within either vertex class is bold. The duals of all remaining edges form the single cycle below. Circle and square classes remain distinguishable without colour.}
\end{figure}
\clearpage
\section*{4. state}
\begin{figure}[h!]\centering
\begin{tikzpicture}[x=0.72cm,y=0.72cm]
\path[use as bounding box] (-.6,-.7) rectangle (12.6,11);
\fill[black!9] (0.00000,0.00000) -- (1.32113,1.43991) -- (4.04994,0.42777) -- cycle;
\fill[black!9] (6.00000,10.39200) -- (2.69330,2.96621) -- (1.32113,1.43991) -- cycle;
\fill[black!9] (6.00000,10.39200) -- (6.18899,7.50100) -- (4.51516,5.61611) -- cycle;
\fill[black!9] (6.00000,10.39200) -- (10.01414,2.60858) -- (8.02049,5.26608) -- cycle;
\fill[black!9] (12.00000,0.00000) -- (4.04994,0.42777) -- (9.22826,0.75499) -- cycle;
\fill[black!9] (12.00000,0.00000) -- (7.28809,3.25717) -- (10.01414,2.60858) -- cycle;
\fill[black!9] (4.04994,0.42777) -- (2.69330,2.96621) -- (6.24482,1.48372) -- cycle;
\fill[black!9] (9.22826,0.75499) -- (6.24482,1.48372) -- (7.28809,3.25717) -- cycle;
\fill[black!9] (2.69330,2.96621) -- (4.51516,5.61611) -- (5.23909,3.48455) -- cycle;
\fill[black!9] (8.02049,5.26608) -- (7.28809,3.25717) -- (6.18899,7.50100) -- cycle;
\fill[black!9] (6.18899,7.50100) -- (6.24482,1.48372) -- (5.23909,3.48455) -- cycle;
\draw[base] (0.00000,0.00000) -- (6.00000,10.39200);
\draw[base] (0.00000,0.00000) -- (12.00000,0.00000);
\draw[base] (0.00000,0.00000) -- (4.04994,0.42777);
\draw[base] (0.00000,0.00000) -- (1.32113,1.43991);
\draw[base] (6.00000,10.39200) -- (12.00000,0.00000);
\draw[base] (6.00000,10.39200) -- (1.32113,1.43991);
\draw[base] (6.00000,10.39200) -- (10.01414,2.60858);
\draw[base] (6.00000,10.39200) -- (2.69330,2.96621);
\draw[base] (6.00000,10.39200) -- (8.02049,5.26608);
\draw[base] (6.00000,10.39200) -- (4.51516,5.61611);
\draw[base] (6.00000,10.39200) -- (6.18899,7.50100);
\draw[base] (12.00000,0.00000) -- (4.04994,0.42777);
\draw[base] (12.00000,0.00000) -- (10.01414,2.60858);
\draw[base] (12.00000,0.00000) -- (9.22826,0.75499);
\draw[base] (12.00000,0.00000) -- (7.28809,3.25717);
\draw[base] (4.04994,0.42777) -- (1.32113,1.43991);
\draw[base] (4.04994,0.42777) -- (9.22826,0.75499);
\draw[base] (4.04994,0.42777) -- (2.69330,2.96621);
\draw[base] (4.04994,0.42777) -- (6.24482,1.48372);
\draw[base] (1.32113,1.43991) -- (2.69330,2.96621);
\draw[base] (10.01414,2.60858) -- (8.02049,5.26608);
\draw[base] (10.01414,2.60858) -- (7.28809,3.25717);
\draw[base] (9.22826,0.75499) -- (7.28809,3.25717);
\draw[base] (9.22826,0.75499) -- (6.24482,1.48372);
\draw[base] (2.69330,2.96621) -- (4.51516,5.61611);
\draw[base] (2.69330,2.96621) -- (5.23909,3.48455);
\draw[base] (2.69330,2.96621) -- (6.24482,1.48372);
\draw[base] (8.02049,5.26608) -- (7.28809,3.25717);
\draw[base] (8.02049,5.26608) -- (6.18899,7.50100);
\draw[base] (7.28809,3.25717) -- (6.18899,7.50100);
\draw[base] (7.28809,3.25717) -- (6.24482,1.48372);
\draw[base] (4.51516,5.61611) -- (6.18899,7.50100);
\draw[base] (4.51516,5.61611) -- (5.23909,3.48455);
\draw[base] (6.18899,7.50100) -- (5.23909,3.48455);
\draw[base] (6.18899,7.50100) -- (6.24482,1.48372);
\draw[base] (5.23909,3.48455) -- (6.24482,1.48372);
\node[markS] at (3.50764,0.47452) {};
\node[markS] at (1.80521,2.27818) {};
\node[markS] at (4.76786,6.14897) {};
\node[markS] at (8.01834,5.46356) {};
\node[markS] at (9.03574,0.66841) {};
\node[markS] at (9.95492,2.45178) {};
\node[markS] at (3.08595,2.64453) {};
\node[markS] at (6.56696,1.56730) {};
\node[markS] at (4.97751,3.61361) {};
\node[markS] at (7.25875,3.75739) {};
\node[markS] at (6.11746,6.69830) {};
\fill (0.00000,0.00000) circle[radius=1.1pt];
\draw[line width=.7pt,fill=white] (0.00000,0.00000) circle[radius=2.4pt];
\fill (12.00000,0.00000) circle[radius=1.1pt];
\draw[line width=.7pt,fill=white] (12.00000,0.00000) circle[radius=2.4pt];
\fill (6.00000,10.39200) circle[radius=1.1pt];
\draw[line width=.7pt,fill=white] (6.00000,10.39200) circle[radius=2.4pt];
\fill (4.04994,0.42777) circle[radius=1.1pt];
\fill (1.32113,1.43991) circle[radius=1.1pt];
\fill (10.01414,2.60858) circle[radius=1.1pt];
\fill (9.22826,0.75499) circle[radius=1.1pt];
\fill (2.69330,2.96621) circle[radius=1.1pt];
\fill (8.02049,5.26608) circle[radius=1.1pt];
\fill (7.28809,3.25717) circle[radius=1.1pt];
\fill (4.51516,5.61611) circle[radius=1.1pt];
\fill (6.18899,7.50100) circle[radius=1.1pt];
\fill (5.23909,3.48455) circle[radius=1.1pt];
\fill (6.24482,1.48372) circle[radius=1.1pt];
\node[anchor=north] at (6,-.25) {$t_0$ is the outer black face};
\end{tikzpicture}
\caption{A blue disk marks the corner $s(t)$ inside each bounded black triangle. Every nonroot receives exactly one mark; roots receive none. This state has a disjoint partner and a forest $Q_s$.}
\end{figure}
\clearpage
\section*{5. pair D}
\begin{figure}[h!]\centering
\begin{tabular}{@{}c@{}}
\begin{tikzpicture}[x=0.72cm,y=0.72cm]
\path[use as bounding box] (-.6,-.7) rectangle (12.6,11);
\fill[black!9] (0.00000,0.00000) -- (1.32113,1.43991) -- (4.04994,0.42777) -- cycle;
\fill[black!9] (6.00000,10.39200) -- (2.69330,2.96621) -- (1.32113,1.43991) -- cycle;
\fill[black!9] (6.00000,10.39200) -- (6.18899,7.50100) -- (4.51516,5.61611) -- cycle;
\fill[black!9] (6.00000,10.39200) -- (10.01414,2.60858) -- (8.02049,5.26608) -- cycle;
\fill[black!9] (12.00000,0.00000) -- (4.04994,0.42777) -- (9.22826,0.75499) -- cycle;
\fill[black!9] (12.00000,0.00000) -- (7.28809,3.25717) -- (10.01414,2.60858) -- cycle;
\fill[black!9] (4.04994,0.42777) -- (2.69330,2.96621) -- (6.24482,1.48372) -- cycle;
\fill[black!9] (9.22826,0.75499) -- (6.24482,1.48372) -- (7.28809,3.25717) -- cycle;
\fill[black!9] (2.69330,2.96621) -- (4.51516,5.61611) -- (5.23909,3.48455) -- cycle;
\fill[black!9] (8.02049,5.26608) -- (7.28809,3.25717) -- (6.18899,7.50100) -- cycle;
\fill[black!9] (6.18899,7.50100) -- (6.24482,1.48372) -- (5.23909,3.48455) -- cycle;
\draw[base] (0.00000,0.00000) -- (6.00000,10.39200);
\draw[base] (0.00000,0.00000) -- (12.00000,0.00000);
\draw[base] (0.00000,0.00000) -- (4.04994,0.42777);
\draw[base] (0.00000,0.00000) -- (1.32113,1.43991);
\draw[base] (6.00000,10.39200) -- (12.00000,0.00000);
\draw[base] (6.00000,10.39200) -- (1.32113,1.43991);
\draw[base] (6.00000,10.39200) -- (10.01414,2.60858);
\draw[base] (6.00000,10.39200) -- (2.69330,2.96621);
\draw[base] (6.00000,10.39200) -- (8.02049,5.26608);
\draw[base] (6.00000,10.39200) -- (4.51516,5.61611);
\draw[base] (6.00000,10.39200) -- (6.18899,7.50100);
\draw[base] (12.00000,0.00000) -- (4.04994,0.42777);
\draw[base] (12.00000,0.00000) -- (10.01414,2.60858);
\draw[base] (12.00000,0.00000) -- (9.22826,0.75499);
\draw[base] (12.00000,0.00000) -- (7.28809,3.25717);
\draw[base] (4.04994,0.42777) -- (1.32113,1.43991);
\draw[base] (4.04994,0.42777) -- (9.22826,0.75499);
\draw[base] (4.04994,0.42777) -- (2.69330,2.96621);
\draw[base] (4.04994,0.42777) -- (6.24482,1.48372);
\draw[base] (1.32113,1.43991) -- (2.69330,2.96621);
\draw[base] (10.01414,2.60858) -- (8.02049,5.26608);
\draw[base] (10.01414,2.60858) -- (7.28809,3.25717);
\draw[base] (9.22826,0.75499) -- (7.28809,3.25717);
\draw[base] (9.22826,0.75499) -- (6.24482,1.48372);
\draw[base] (2.69330,2.96621) -- (4.51516,5.61611);
\draw[base] (2.69330,2.96621) -- (5.23909,3.48455);
\draw[base] (2.69330,2.96621) -- (6.24482,1.48372);
\draw[base] (8.02049,5.26608) -- (7.28809,3.25717);
\draw[base] (8.02049,5.26608) -- (6.18899,7.50100);
\draw[base] (7.28809,3.25717) -- (6.18899,7.50100);
\draw[base] (7.28809,3.25717) -- (6.24482,1.48372);
\draw[base] (4.51516,5.61611) -- (6.18899,7.50100);
\draw[base] (4.51516,5.61611) -- (5.23909,3.48455);
\draw[base] (6.18899,7.50100) -- (5.23909,3.48455);
\draw[base] (6.18899,7.50100) -- (6.24482,1.48372);
\draw[base] (5.23909,3.48455) -- (6.24482,1.48372);
\node[markS] at (1.43374,1.24375) {};
\node[markR] at (3.50764,0.47452) {};
\node[markS] at (2.84806,3.43817) {};
\node[markR] at (1.80521,2.27818) {};
\node[markS] at (6.03996,7.58149) {};
\node[markR] at (4.76786,6.14897) {};
\node[markS] at (9.53351,3.44385) {};
\node[markR] at (8.01834,5.46356) {};
\node[markS] at (5.10021,0.41973) {};
\node[markR] at (9.03574,0.66841) {};
\node[markS] at (7.88313,2.94471) {};
\node[markR] at (9.95492,2.45178) {};
\node[markS] at (5.78511,1.51785) {};
\node[markR] at (3.08595,2.64453) {};
\node[markS] at (8.83437,1.01346) {};
\node[markR] at (6.56696,1.56730) {};
\node[markS] at (4.42733,5.23359) {};
\node[markR] at (4.97751,3.61361) {};
\node[markS] at (7.81538,5.28416) {};
\node[markR] at (7.25875,3.75739) {};
\node[markS] at (5.39554,3.64580) {};
\node[markR] at (6.11746,6.69830) {};
\draw[dedge] (4.04994,0.42777) -- (1.32113,1.43991);
\draw[dedge] (1.32113,1.43991) -- (2.69330,2.96621);
\draw[dedge] (4.51516,5.61611) -- (6.18899,7.50100);
\draw[dedge] (8.02049,5.26608) -- (10.01414,2.60858);
\draw[dedge] (9.22826,0.75499) -- (4.04994,0.42777);
\draw[dedge] (10.01414,2.60858) -- (7.28809,3.25717);
\draw[dedge] (2.69330,2.96621) -- (6.24482,1.48372);
\draw[dedge] (6.24482,1.48372) -- (9.22826,0.75499);
\draw[dedge] (5.23909,3.48455) -- (4.51516,5.61611);
\draw[dedge] (7.28809,3.25717) -- (8.02049,5.26608);
\draw[dedge] (6.18899,7.50100) -- (5.23909,3.48455);
\fill (0.00000,0.00000) circle[radius=1.1pt];
\draw[line width=.7pt,fill=white] (0.00000,0.00000) circle[radius=2.4pt];
\fill (12.00000,0.00000) circle[radius=1.1pt];
\draw[line width=.7pt,fill=white] (12.00000,0.00000) circle[radius=2.4pt];
\fill (6.00000,10.39200) circle[radius=1.1pt];
\draw[line width=.7pt,fill=white] (6.00000,10.39200) circle[radius=2.4pt];
\fill (4.04994,0.42777) circle[radius=1.1pt];
\fill (1.32113,1.43991) circle[radius=1.1pt];
\fill (10.01414,2.60858) circle[radius=1.1pt];
\fill (9.22826,0.75499) circle[radius=1.1pt];
\fill (2.69330,2.96621) circle[radius=1.1pt];
\fill (8.02049,5.26608) circle[radius=1.1pt];
\fill (7.28809,3.25717) circle[radius=1.1pt];
\fill (4.51516,5.61611) circle[radius=1.1pt];
\fill (6.18899,7.50100) circle[radius=1.1pt];
\fill (5.23909,3.48455) circle[radius=1.1pt];
\fill (6.24482,1.48372) circle[radius=1.1pt];
\node[anchor=north] at (6,-.25) {$t_0$ is the outer black face};
\end{tikzpicture}\\[1mm]Squares: $r$; disks: $s$\\[4mm]\begin{tikzpicture}[x=0.72cm,y=0.72cm]
\path[use as bounding box] (-.6,-.7) rectangle (12.6,11);
\fill[black!9] (0.00000,0.00000) -- (1.32113,1.43991) -- (4.04994,0.42777) -- cycle;
\fill[black!9] (6.00000,10.39200) -- (2.69330,2.96621) -- (1.32113,1.43991) -- cycle;
\fill[black!9] (6.00000,10.39200) -- (6.18899,7.50100) -- (4.51516,5.61611) -- cycle;
\fill[black!9] (6.00000,10.39200) -- (10.01414,2.60858) -- (8.02049,5.26608) -- cycle;
\fill[black!9] (12.00000,0.00000) -- (4.04994,0.42777) -- (9.22826,0.75499) -- cycle;
\fill[black!9] (12.00000,0.00000) -- (7.28809,3.25717) -- (10.01414,2.60858) -- cycle;
\fill[black!9] (4.04994,0.42777) -- (2.69330,2.96621) -- (6.24482,1.48372) -- cycle;
\fill[black!9] (9.22826,0.75499) -- (6.24482,1.48372) -- (7.28809,3.25717) -- cycle;
\fill[black!9] (2.69330,2.96621) -- (4.51516,5.61611) -- (5.23909,3.48455) -- cycle;
\fill[black!9] (8.02049,5.26608) -- (7.28809,3.25717) -- (6.18899,7.50100) -- cycle;
\fill[black!9] (6.18899,7.50100) -- (6.24482,1.48372) -- (5.23909,3.48455) -- cycle;
\draw[base] (0.00000,0.00000) -- (6.00000,10.39200);
\draw[base] (0.00000,0.00000) -- (12.00000,0.00000);
\draw[base] (0.00000,0.00000) -- (4.04994,0.42777);
\draw[base] (0.00000,0.00000) -- (1.32113,1.43991);
\draw[base] (6.00000,10.39200) -- (12.00000,0.00000);
\draw[base] (6.00000,10.39200) -- (1.32113,1.43991);
\draw[base] (6.00000,10.39200) -- (10.01414,2.60858);
\draw[base] (6.00000,10.39200) -- (2.69330,2.96621);
\draw[base] (6.00000,10.39200) -- (8.02049,5.26608);
\draw[base] (6.00000,10.39200) -- (4.51516,5.61611);
\draw[base] (6.00000,10.39200) -- (6.18899,7.50100);
\draw[base] (12.00000,0.00000) -- (4.04994,0.42777);
\draw[base] (12.00000,0.00000) -- (10.01414,2.60858);
\draw[base] (12.00000,0.00000) -- (9.22826,0.75499);
\draw[base] (12.00000,0.00000) -- (7.28809,3.25717);
\draw[base] (4.04994,0.42777) -- (1.32113,1.43991);
\draw[base] (4.04994,0.42777) -- (9.22826,0.75499);
\draw[base] (4.04994,0.42777) -- (2.69330,2.96621);
\draw[base] (4.04994,0.42777) -- (6.24482,1.48372);
\draw[base] (1.32113,1.43991) -- (2.69330,2.96621);
\draw[base] (10.01414,2.60858) -- (8.02049,5.26608);
\draw[base] (10.01414,2.60858) -- (7.28809,3.25717);
\draw[base] (9.22826,0.75499) -- (7.28809,3.25717);
\draw[base] (9.22826,0.75499) -- (6.24482,1.48372);
\draw[base] (2.69330,2.96621) -- (4.51516,5.61611);
\draw[base] (2.69330,2.96621) -- (5.23909,3.48455);
\draw[base] (2.69330,2.96621) -- (6.24482,1.48372);
\draw[base] (8.02049,5.26608) -- (7.28809,3.25717);
\draw[base] (8.02049,5.26608) -- (6.18899,7.50100);
\draw[base] (7.28809,3.25717) -- (6.18899,7.50100);
\draw[base] (7.28809,3.25717) -- (6.24482,1.48372);
\draw[base] (4.51516,5.61611) -- (6.18899,7.50100);
\draw[base] (4.51516,5.61611) -- (5.23909,3.48455);
\draw[base] (6.18899,7.50100) -- (5.23909,3.48455);
\draw[base] (6.18899,7.50100) -- (6.24482,1.48372);
\draw[base] (5.23909,3.48455) -- (6.24482,1.48372);
\draw[dedge] (4.04994,0.42777) -- (1.32113,1.43991);
\draw[dedge] (1.32113,1.43991) -- (2.69330,2.96621);
\draw[dedge] (4.51516,5.61611) -- (6.18899,7.50100);
\draw[dedge] (8.02049,5.26608) -- (10.01414,2.60858);
\draw[dedge] (9.22826,0.75499) -- (4.04994,0.42777);
\draw[dedge] (10.01414,2.60858) -- (7.28809,3.25717);
\draw[dedge] (2.69330,2.96621) -- (6.24482,1.48372);
\draw[dedge] (6.24482,1.48372) -- (9.22826,0.75499);
\draw[dedge] (5.23909,3.48455) -- (4.51516,5.61611);
\draw[dedge] (7.28809,3.25717) -- (8.02049,5.26608);
\draw[dedge] (6.18899,7.50100) -- (5.23909,3.48455);
\fill (0.00000,0.00000) circle[radius=1.1pt];
\draw[line width=.7pt,fill=white] (0.00000,0.00000) circle[radius=2.4pt];
\fill (12.00000,0.00000) circle[radius=1.1pt];
\draw[line width=.7pt,fill=white] (12.00000,0.00000) circle[radius=2.4pt];
\fill (6.00000,10.39200) circle[radius=1.1pt];
\draw[line width=.7pt,fill=white] (6.00000,10.39200) circle[radius=2.4pt];
\fill (4.04994,0.42777) circle[radius=1.1pt];
\fill (1.32113,1.43991) circle[radius=1.1pt];
\fill (10.01414,2.60858) circle[radius=1.1pt];
\fill (9.22826,0.75499) circle[radius=1.1pt];
\fill (2.69330,2.96621) circle[radius=1.1pt];
\fill (8.02049,5.26608) circle[radius=1.1pt];
\fill (7.28809,3.25717) circle[radius=1.1pt];
\fill (4.51516,5.61611) circle[radius=1.1pt];
\fill (6.18899,7.50100) circle[radius=1.1pt];
\fill (5.23909,3.48455) circle[radius=1.1pt];
\fill (6.24482,1.48372) circle[radius=1.1pt];
\node[anchor=north] at (6,-.25) {$t_0$ is the outer black face};
\end{tikzpicture}\\[1mm]$D$: arrows from $r(t)$ to $s(t)$
\end{tabular}
\caption{The two states choose different corners in every black face. Their directed edges split the nonroots into three cycles. The lower panel removes corner marks to expose those cycles.}
\end{figure}
\clearpage
\section*{6. Q}
\begin{figure}[h!]\centering
\begin{tabular}{@{}c@{}}
\begin{tikzpicture}[x=0.72cm,y=0.72cm]
\path[use as bounding box] (-.6,-.7) rectangle (12.6,11);
\fill[black!9] (0.00000,0.00000) -- (1.32113,1.43991) -- (4.04994,0.42777) -- cycle;
\fill[black!9] (6.00000,10.39200) -- (2.69330,2.96621) -- (1.32113,1.43991) -- cycle;
\fill[black!9] (6.00000,10.39200) -- (6.18899,7.50100) -- (4.51516,5.61611) -- cycle;
\fill[black!9] (6.00000,10.39200) -- (10.01414,2.60858) -- (8.02049,5.26608) -- cycle;
\fill[black!9] (12.00000,0.00000) -- (4.04994,0.42777) -- (9.22826,0.75499) -- cycle;
\fill[black!9] (12.00000,0.00000) -- (7.28809,3.25717) -- (10.01414,2.60858) -- cycle;
\fill[black!9] (4.04994,0.42777) -- (2.69330,2.96621) -- (6.24482,1.48372) -- cycle;
\fill[black!9] (9.22826,0.75499) -- (6.24482,1.48372) -- (7.28809,3.25717) -- cycle;
\fill[black!9] (2.69330,2.96621) -- (4.51516,5.61611) -- (5.23909,3.48455) -- cycle;
\fill[black!9] (8.02049,5.26608) -- (7.28809,3.25717) -- (6.18899,7.50100) -- cycle;
\fill[black!9] (6.18899,7.50100) -- (6.24482,1.48372) -- (5.23909,3.48455) -- cycle;
\draw[base] (0.00000,0.00000) -- (6.00000,10.39200);
\draw[base] (0.00000,0.00000) -- (12.00000,0.00000);
\draw[base] (0.00000,0.00000) -- (4.04994,0.42777);
\draw[base] (0.00000,0.00000) -- (1.32113,1.43991);
\draw[base] (6.00000,10.39200) -- (12.00000,0.00000);
\draw[base] (6.00000,10.39200) -- (1.32113,1.43991);
\draw[base] (6.00000,10.39200) -- (10.01414,2.60858);
\draw[base] (6.00000,10.39200) -- (2.69330,2.96621);
\draw[base] (6.00000,10.39200) -- (8.02049,5.26608);
\draw[base] (6.00000,10.39200) -- (4.51516,5.61611);
\draw[base] (6.00000,10.39200) -- (6.18899,7.50100);
\draw[base] (12.00000,0.00000) -- (4.04994,0.42777);
\draw[base] (12.00000,0.00000) -- (10.01414,2.60858);
\draw[base] (12.00000,0.00000) -- (9.22826,0.75499);
\draw[base] (12.00000,0.00000) -- (7.28809,3.25717);
\draw[base] (4.04994,0.42777) -- (1.32113,1.43991);
\draw[base] (4.04994,0.42777) -- (9.22826,0.75499);
\draw[base] (4.04994,0.42777) -- (2.69330,2.96621);
\draw[base] (4.04994,0.42777) -- (6.24482,1.48372);
\draw[base] (1.32113,1.43991) -- (2.69330,2.96621);
\draw[base] (10.01414,2.60858) -- (8.02049,5.26608);
\draw[base] (10.01414,2.60858) -- (7.28809,3.25717);
\draw[base] (9.22826,0.75499) -- (7.28809,3.25717);
\draw[base] (9.22826,0.75499) -- (6.24482,1.48372);
\draw[base] (2.69330,2.96621) -- (4.51516,5.61611);
\draw[base] (2.69330,2.96621) -- (5.23909,3.48455);
\draw[base] (2.69330,2.96621) -- (6.24482,1.48372);
\draw[base] (8.02049,5.26608) -- (7.28809,3.25717);
\draw[base] (8.02049,5.26608) -- (6.18899,7.50100);
\draw[base] (7.28809,3.25717) -- (6.18899,7.50100);
\draw[base] (7.28809,3.25717) -- (6.24482,1.48372);
\draw[base] (4.51516,5.61611) -- (6.18899,7.50100);
\draw[base] (4.51516,5.61611) -- (5.23909,3.48455);
\draw[base] (6.18899,7.50100) -- (5.23909,3.48455);
\draw[base] (6.18899,7.50100) -- (6.24482,1.48372);
\draw[base] (5.23909,3.48455) -- (6.24482,1.48372);
\draw[qedge,-{Stealth[length=3pt]}] (4.04994,0.42777) -- (0.00000,0.00000);
\draw[qedge,-{Stealth[length=3pt]}] (1.32113,1.43991) -- (6.00000,10.39200);
\draw[qedge,-{Stealth[length=3pt]}] (4.51516,5.61611) -- (6.00000,10.39200);
\draw[qedge,-{Stealth[length=3pt]}] (8.02049,5.26608) -- (6.00000,10.39200);
\draw[qedge,-{Stealth[length=3pt]}] (9.22826,0.75499) -- (12.00000,0.00000);
\draw[qedge,-{Stealth[length=3pt]}] (10.01414,2.60858) -- (12.00000,0.00000);
\draw[qedge,-{Stealth[length=3pt]}] (2.69330,2.96621) -- (4.04994,0.42777);
\draw[qedge,-{Stealth[length=3pt]}] (6.24482,1.48372) -- (7.28809,3.25717);
\draw[qedge,-{Stealth[length=3pt]}] (5.23909,3.48455) -- (2.69330,2.96621);
\draw[qedge,-{Stealth[length=3pt]}] (7.28809,3.25717) -- (6.18899,7.50100);
\draw[qedge,-{Stealth[length=3pt]}] (6.18899,7.50100) -- (6.24482,1.48372);
\fill (0.00000,0.00000) circle[radius=1.1pt];
\draw[line width=.7pt,fill=white] (0.00000,0.00000) circle[radius=2.4pt];
\fill (12.00000,0.00000) circle[radius=1.1pt];
\draw[line width=.7pt,fill=white] (12.00000,0.00000) circle[radius=2.4pt];
\fill (6.00000,10.39200) circle[radius=1.1pt];
\draw[line width=.7pt,fill=white] (6.00000,10.39200) circle[radius=2.4pt];
\fill (4.04994,0.42777) circle[radius=1.1pt];
\fill (1.32113,1.43991) circle[radius=1.1pt];
\fill (10.01414,2.60858) circle[radius=1.1pt];
\fill (9.22826,0.75499) circle[radius=1.1pt];
\fill (2.69330,2.96621) circle[radius=1.1pt];
\fill (8.02049,5.26608) circle[radius=1.1pt];
\fill (7.28809,3.25717) circle[radius=1.1pt];
\fill (4.51516,5.61611) circle[radius=1.1pt];
\fill (6.18899,7.50100) circle[radius=1.1pt];
\fill (5.23909,3.48455) circle[radius=1.1pt];
\fill (6.24482,1.48372) circle[radius=1.1pt];
\node[anchor=north] at (6,-.25) {$t_0$ is the outer black face};
\end{tikzpicture}\\[1mm]Cyclic $Q_s$\\[4mm]\begin{tikzpicture}[x=0.72cm,y=0.72cm]
\path[use as bounding box] (-.6,-.7) rectangle (12.6,11);
\fill[black!9] (0.00000,0.00000) -- (1.32113,1.43991) -- (4.04994,0.42777) -- cycle;
\fill[black!9] (6.00000,10.39200) -- (2.69330,2.96621) -- (1.32113,1.43991) -- cycle;
\fill[black!9] (6.00000,10.39200) -- (6.18899,7.50100) -- (4.51516,5.61611) -- cycle;
\fill[black!9] (6.00000,10.39200) -- (10.01414,2.60858) -- (8.02049,5.26608) -- cycle;
\fill[black!9] (12.00000,0.00000) -- (4.04994,0.42777) -- (9.22826,0.75499) -- cycle;
\fill[black!9] (12.00000,0.00000) -- (7.28809,3.25717) -- (10.01414,2.60858) -- cycle;
\fill[black!9] (4.04994,0.42777) -- (2.69330,2.96621) -- (6.24482,1.48372) -- cycle;
\fill[black!9] (9.22826,0.75499) -- (6.24482,1.48372) -- (7.28809,3.25717) -- cycle;
\fill[black!9] (2.69330,2.96621) -- (4.51516,5.61611) -- (5.23909,3.48455) -- cycle;
\fill[black!9] (8.02049,5.26608) -- (7.28809,3.25717) -- (6.18899,7.50100) -- cycle;
\fill[black!9] (6.18899,7.50100) -- (6.24482,1.48372) -- (5.23909,3.48455) -- cycle;
\draw[base] (0.00000,0.00000) -- (6.00000,10.39200);
\draw[base] (0.00000,0.00000) -- (12.00000,0.00000);
\draw[base] (0.00000,0.00000) -- (4.04994,0.42777);
\draw[base] (0.00000,0.00000) -- (1.32113,1.43991);
\draw[base] (6.00000,10.39200) -- (12.00000,0.00000);
\draw[base] (6.00000,10.39200) -- (1.32113,1.43991);
\draw[base] (6.00000,10.39200) -- (10.01414,2.60858);
\draw[base] (6.00000,10.39200) -- (2.69330,2.96621);
\draw[base] (6.00000,10.39200) -- (8.02049,5.26608);
\draw[base] (6.00000,10.39200) -- (4.51516,5.61611);
\draw[base] (6.00000,10.39200) -- (6.18899,7.50100);
\draw[base] (12.00000,0.00000) -- (4.04994,0.42777);
\draw[base] (12.00000,0.00000) -- (10.01414,2.60858);
\draw[base] (12.00000,0.00000) -- (9.22826,0.75499);
\draw[base] (12.00000,0.00000) -- (7.28809,3.25717);
\draw[base] (4.04994,0.42777) -- (1.32113,1.43991);
\draw[base] (4.04994,0.42777) -- (9.22826,0.75499);
\draw[base] (4.04994,0.42777) -- (2.69330,2.96621);
\draw[base] (4.04994,0.42777) -- (6.24482,1.48372);
\draw[base] (1.32113,1.43991) -- (2.69330,2.96621);
\draw[base] (10.01414,2.60858) -- (8.02049,5.26608);
\draw[base] (10.01414,2.60858) -- (7.28809,3.25717);
\draw[base] (9.22826,0.75499) -- (7.28809,3.25717);
\draw[base] (9.22826,0.75499) -- (6.24482,1.48372);
\draw[base] (2.69330,2.96621) -- (4.51516,5.61611);
\draw[base] (2.69330,2.96621) -- (5.23909,3.48455);
\draw[base] (2.69330,2.96621) -- (6.24482,1.48372);
\draw[base] (8.02049,5.26608) -- (7.28809,3.25717);
\draw[base] (8.02049,5.26608) -- (6.18899,7.50100);
\draw[base] (7.28809,3.25717) -- (6.18899,7.50100);
\draw[base] (7.28809,3.25717) -- (6.24482,1.48372);
\draw[base] (4.51516,5.61611) -- (6.18899,7.50100);
\draw[base] (4.51516,5.61611) -- (5.23909,3.48455);
\draw[base] (6.18899,7.50100) -- (5.23909,3.48455);
\draw[base] (6.18899,7.50100) -- (6.24482,1.48372);
\draw[base] (5.23909,3.48455) -- (6.24482,1.48372);
\draw[qedge,-{Stealth[length=3pt]}] (1.32113,1.43991) -- (0.00000,0.00000);
\draw[qedge,-{Stealth[length=3pt]}] (2.69330,2.96621) -- (6.00000,10.39200);
\draw[qedge,-{Stealth[length=3pt]}] (6.18899,7.50100) -- (6.00000,10.39200);
\draw[qedge,-{Stealth[length=3pt]}] (10.01414,2.60858) -- (6.00000,10.39200);
\draw[qedge,-{Stealth[length=3pt]}] (4.04994,0.42777) -- (12.00000,0.00000);
\draw[qedge,-{Stealth[length=3pt]}] (7.28809,3.25717) -- (12.00000,0.00000);
\draw[qedge,-{Stealth[length=3pt]}] (6.24482,1.48372) -- (4.04994,0.42777);
\draw[qedge,-{Stealth[length=3pt]}] (9.22826,0.75499) -- (7.28809,3.25717);
\draw[qedge,-{Stealth[length=3pt]}] (4.51516,5.61611) -- (2.69330,2.96621);
\draw[qedge,-{Stealth[length=3pt]}] (8.02049,5.26608) -- (6.18899,7.50100);
\draw[qedge,-{Stealth[length=3pt]}] (5.23909,3.48455) -- (6.24482,1.48372);
\fill (0.00000,0.00000) circle[radius=1.1pt];
\draw[line width=.7pt,fill=white] (0.00000,0.00000) circle[radius=2.4pt];
\fill (12.00000,0.00000) circle[radius=1.1pt];
\draw[line width=.7pt,fill=white] (12.00000,0.00000) circle[radius=2.4pt];
\fill (6.00000,10.39200) circle[radius=1.1pt];
\draw[line width=.7pt,fill=white] (6.00000,10.39200) circle[radius=2.4pt];
\fill (4.04994,0.42777) circle[radius=1.1pt];
\fill (1.32113,1.43991) circle[radius=1.1pt];
\fill (10.01414,2.60858) circle[radius=1.1pt];
\fill (9.22826,0.75499) circle[radius=1.1pt];
\fill (2.69330,2.96621) circle[radius=1.1pt];
\fill (8.02049,5.26608) circle[radius=1.1pt];
\fill (7.28809,3.25717) circle[radius=1.1pt];
\fill (4.51516,5.61611) circle[radius=1.1pt];
\fill (6.18899,7.50100) circle[radius=1.1pt];
\fill (5.23909,3.48455) circle[radius=1.1pt];
\fill (6.24482,1.48372) circle[radius=1.1pt];
\node[anchor=north] at (6,-.25) {$t_0$ is the outer black face};
\end{tikzpicture}\\[1mm]Forest $Q_s$
\end{tabular}
\caption{Each dashed edge lies opposite the second-state mark. The upper state has a cycle; the lower has three tree components, each containing one root. The panels use different second states.}
\end{figure}
\clearpage
\section*{7. disk}
\begin{figure}[h!]\centering
\begin{tikzpicture}[x=0.72cm,y=0.72cm]
\path[use as bounding box] (-1.5,-1.4) rectangle (13.5,11);
\fill[black!9] (0.00000,0.00000) -- (1.32113,1.43991) -- (4.04994,0.42777) -- cycle;
\fill[black!9] (6.00000,10.39200) -- (2.69330,2.96621) -- (1.32113,1.43991) -- cycle;
\fill[black!9] (6.00000,10.39200) -- (6.18899,7.50100) -- (4.51516,5.61611) -- cycle;
\fill[black!9] (6.00000,10.39200) -- (10.01414,2.60858) -- (8.02049,5.26608) -- cycle;
\fill[black!9] (12.00000,0.00000) -- (4.04994,0.42777) -- (9.22826,0.75499) -- cycle;
\fill[black!9] (12.00000,0.00000) -- (7.28809,3.25717) -- (10.01414,2.60858) -- cycle;
\fill[black!9] (4.04994,0.42777) -- (2.69330,2.96621) -- (6.24482,1.48372) -- cycle;
\fill[black!9] (9.22826,0.75499) -- (6.24482,1.48372) -- (7.28809,3.25717) -- cycle;
\fill[black!9] (2.69330,2.96621) -- (4.51516,5.61611) -- (5.23909,3.48455) -- cycle;
\fill[black!9] (8.02049,5.26608) -- (7.28809,3.25717) -- (6.18899,7.50100) -- cycle;
\fill[black!9] (6.18899,7.50100) -- (6.24482,1.48372) -- (5.23909,3.48455) -- cycle;
\fill[pairorange!18] (4.04994,0.42777) -- (1.32113,1.43991) -- (2.69330,2.96621) -- cycle;
\fill[pairorange!18] (4.04994,0.42777) -- (6.24482,1.48372) -- (9.22826,0.75499) -- cycle;
\fill[pairorange!18] (4.04994,0.42777) -- (2.69330,2.96621) -- (6.24482,1.48372) -- cycle;
\draw[base] (0.00000,0.00000) -- (6.00000,10.39200);
\draw[base] (0.00000,0.00000) -- (12.00000,0.00000);
\draw[base] (0.00000,0.00000) -- (4.04994,0.42777);
\draw[base] (0.00000,0.00000) -- (1.32113,1.43991);
\draw[base] (6.00000,10.39200) -- (12.00000,0.00000);
\draw[base] (6.00000,10.39200) -- (1.32113,1.43991);
\draw[base] (6.00000,10.39200) -- (10.01414,2.60858);
\draw[base] (6.00000,10.39200) -- (2.69330,2.96621);
\draw[base] (6.00000,10.39200) -- (8.02049,5.26608);
\draw[base] (6.00000,10.39200) -- (4.51516,5.61611);
\draw[base] (6.00000,10.39200) -- (6.18899,7.50100);
\draw[base] (12.00000,0.00000) -- (4.04994,0.42777);
\draw[base] (12.00000,0.00000) -- (10.01414,2.60858);
\draw[base] (12.00000,0.00000) -- (9.22826,0.75499);
\draw[base] (12.00000,0.00000) -- (7.28809,3.25717);
\draw[base] (4.04994,0.42777) -- (1.32113,1.43991);
\draw[base] (4.04994,0.42777) -- (9.22826,0.75499);
\draw[base] (4.04994,0.42777) -- (2.69330,2.96621);
\draw[base] (4.04994,0.42777) -- (6.24482,1.48372);
\draw[base] (1.32113,1.43991) -- (2.69330,2.96621);
\draw[base] (10.01414,2.60858) -- (8.02049,5.26608);
\draw[base] (10.01414,2.60858) -- (7.28809,3.25717);
\draw[base] (9.22826,0.75499) -- (7.28809,3.25717);
\draw[base] (9.22826,0.75499) -- (6.24482,1.48372);
\draw[base] (2.69330,2.96621) -- (4.51516,5.61611);
\draw[base] (2.69330,2.96621) -- (5.23909,3.48455);
\draw[base] (2.69330,2.96621) -- (6.24482,1.48372);
\draw[base] (8.02049,5.26608) -- (7.28809,3.25717);
\draw[base] (8.02049,5.26608) -- (6.18899,7.50100);
\draw[base] (7.28809,3.25717) -- (6.18899,7.50100);
\draw[base] (7.28809,3.25717) -- (6.24482,1.48372);
\draw[base] (4.51516,5.61611) -- (6.18899,7.50100);
\draw[base] (4.51516,5.61611) -- (5.23909,3.48455);
\draw[base] (6.18899,7.50100) -- (5.23909,3.48455);
\draw[base] (6.18899,7.50100) -- (6.24482,1.48372);
\draw[base] (5.23909,3.48455) -- (6.24482,1.48372);
\draw[dedge] (4.04994,0.42777) -- (1.32113,1.43991);
\draw[dedge] (1.32113,1.43991) -- (2.69330,2.96621);
\draw[dedge] (2.69330,2.96621) -- (6.24482,1.48372);
\draw[dedge] (6.24482,1.48372) -- (9.22826,0.75499);
\draw[dedge] (9.22826,0.75499) -- (4.04994,0.42777);
\fill (0.00000,0.00000) circle[radius=1.1pt];
\draw[line width=.7pt,fill=white] (0.00000,0.00000) circle[radius=2.4pt];
\fill (12.00000,0.00000) circle[radius=1.1pt];
\draw[line width=.7pt,fill=white] (12.00000,0.00000) circle[radius=2.4pt];
\fill (6.00000,10.39200) circle[radius=1.1pt];
\draw[line width=.7pt,fill=white] (6.00000,10.39200) circle[radius=2.4pt];
\fill (4.04994,0.42777) circle[radius=1.1pt];
\fill (1.32113,1.43991) circle[radius=1.1pt];
\fill (10.01414,2.60858) circle[radius=1.1pt];
\fill (9.22826,0.75499) circle[radius=1.1pt];
\fill (2.69330,2.96621) circle[radius=1.1pt];
\fill (8.02049,5.26608) circle[radius=1.1pt];
\fill (7.28809,3.25717) circle[radius=1.1pt];
\fill (4.51516,5.61611) circle[radius=1.1pt];
\fill (6.18899,7.50100) circle[radius=1.1pt];
\fill (5.23909,3.48455) circle[radius=1.1pt];
\fill (6.24482,1.48372) circle[radius=1.1pt];
\node[fill=white,inner sep=1pt] at (2.68812,1.61130) {$W$};
\node[fill=white,inner sep=1pt] at (6.50768,0.88883) {$W$};
\node[fill=white,inner sep=1pt] at (4.32936,1.62590) {$B$};
\node[anchor=north] at (6,-.8) {$B=1,\ W=2,\ I=0,\ h=5,\ l=1$};
\node[anchor=north] at (6,-.25) {$t_0$ is the outer black face};
\end{tikzpicture}
\caption{The highlighted disk bounded by a five-edge cycle has $B=1$, $W=2$, $I=0$, $h=5$ and $l=1$. Thus $B+W=2I+h-2$, $B=I+l$ and $2l-h=-3$. The other two cycles of $D$ are suppressed.}
\end{figure}
\clearpage
\section*{8. cancellation}
\begin{figure}[h!]\centering
\begin{tabular}{@{}c@{}}
\begin{tikzpicture}[x=0.72cm,y=0.72cm]
\path[use as bounding box] (-.6,-.7) rectangle (12.6,11);
\fill[black!9] (0.00000,0.00000) -- (1.32113,1.43991) -- (4.04994,0.42777) -- cycle;
\fill[black!9] (6.00000,10.39200) -- (2.69330,2.96621) -- (1.32113,1.43991) -- cycle;
\fill[black!9] (6.00000,10.39200) -- (6.18899,7.50100) -- (4.51516,5.61611) -- cycle;
\fill[black!9] (6.00000,10.39200) -- (10.01414,2.60858) -- (8.02049,5.26608) -- cycle;
\fill[black!9] (12.00000,0.00000) -- (4.04994,0.42777) -- (9.22826,0.75499) -- cycle;
\fill[black!9] (12.00000,0.00000) -- (7.28809,3.25717) -- (10.01414,2.60858) -- cycle;
\fill[black!9] (4.04994,0.42777) -- (2.69330,2.96621) -- (6.24482,1.48372) -- cycle;
\fill[black!9] (9.22826,0.75499) -- (6.24482,1.48372) -- (7.28809,3.25717) -- cycle;
\fill[black!9] (2.69330,2.96621) -- (4.51516,5.61611) -- (5.23909,3.48455) -- cycle;
\fill[black!9] (8.02049,5.26608) -- (7.28809,3.25717) -- (6.18899,7.50100) -- cycle;
\fill[black!9] (6.18899,7.50100) -- (6.24482,1.48372) -- (5.23909,3.48455) -- cycle;
\draw[base] (0.00000,0.00000) -- (6.00000,10.39200);
\draw[base] (0.00000,0.00000) -- (12.00000,0.00000);
\draw[base] (0.00000,0.00000) -- (4.04994,0.42777);
\draw[base] (0.00000,0.00000) -- (1.32113,1.43991);
\draw[base] (6.00000,10.39200) -- (12.00000,0.00000);
\draw[base] (6.00000,10.39200) -- (1.32113,1.43991);
\draw[base] (6.00000,10.39200) -- (10.01414,2.60858);
\draw[base] (6.00000,10.39200) -- (2.69330,2.96621);
\draw[base] (6.00000,10.39200) -- (8.02049,5.26608);
\draw[base] (6.00000,10.39200) -- (4.51516,5.61611);
\draw[base] (6.00000,10.39200) -- (6.18899,7.50100);
\draw[base] (12.00000,0.00000) -- (4.04994,0.42777);
\draw[base] (12.00000,0.00000) -- (10.01414,2.60858);
\draw[base] (12.00000,0.00000) -- (9.22826,0.75499);
\draw[base] (12.00000,0.00000) -- (7.28809,3.25717);
\draw[base] (4.04994,0.42777) -- (1.32113,1.43991);
\draw[base] (4.04994,0.42777) -- (9.22826,0.75499);
\draw[base] (4.04994,0.42777) -- (2.69330,2.96621);
\draw[base] (4.04994,0.42777) -- (6.24482,1.48372);
\draw[base] (1.32113,1.43991) -- (2.69330,2.96621);
\draw[base] (10.01414,2.60858) -- (8.02049,5.26608);
\draw[base] (10.01414,2.60858) -- (7.28809,3.25717);
\draw[base] (9.22826,0.75499) -- (7.28809,3.25717);
\draw[base] (9.22826,0.75499) -- (6.24482,1.48372);
\draw[base] (2.69330,2.96621) -- (4.51516,5.61611);
\draw[base] (2.69330,2.96621) -- (5.23909,3.48455);
\draw[base] (2.69330,2.96621) -- (6.24482,1.48372);
\draw[base] (8.02049,5.26608) -- (7.28809,3.25717);
\draw[base] (8.02049,5.26608) -- (6.18899,7.50100);
\draw[base] (7.28809,3.25717) -- (6.18899,7.50100);
\draw[base] (7.28809,3.25717) -- (6.24482,1.48372);
\draw[base] (4.51516,5.61611) -- (6.18899,7.50100);
\draw[base] (4.51516,5.61611) -- (5.23909,3.48455);
\draw[base] (6.18899,7.50100) -- (5.23909,3.48455);
\draw[base] (6.18899,7.50100) -- (6.24482,1.48372);
\draw[base] (5.23909,3.48455) -- (6.24482,1.48372);
\node[markS] at (1.43374,1.24375) {};
\node[markR] at (3.50764,0.47452) {};
\node[markS] at (2.84806,3.43817) {};
\node[markR] at (1.80521,2.27818) {};
\node[markS] at (6.03996,7.58149) {};
\node[markR] at (4.76786,6.14897) {};
\node[markS] at (9.53351,3.44385) {};
\node[markR] at (8.01834,5.46356) {};
\node[markS] at (5.10021,0.41973) {};
\node[markR] at (9.03574,0.66841) {};
\node[markS] at (7.88313,2.94471) {};
\node[markR] at (9.95492,2.45178) {};
\node[markS] at (5.78511,1.51785) {};
\node[markR] at (3.08595,2.64453) {};
\node[markS] at (8.83437,1.01346) {};
\node[markR] at (6.56696,1.56730) {};
\node[markS] at (4.42733,5.23359) {};
\node[markR] at (4.97751,3.61361) {};
\node[markS] at (7.81538,5.28416) {};
\node[markR] at (7.25875,3.75739) {};
\node[markS] at (5.39554,3.64580) {};
\node[markR] at (6.11746,6.69830) {};
\draw[qedge,-{Stealth[length=3pt]}] (4.04994,0.42777) -- (0.00000,0.00000);
\draw[qedge,-{Stealth[length=3pt]}] (1.32113,1.43991) -- (6.00000,10.39200);
\draw[qedge,-{Stealth[length=3pt]}] (4.51516,5.61611) -- (6.00000,10.39200);
\draw[qedge,-{Stealth[length=3pt]}] (8.02049,5.26608) -- (6.00000,10.39200);
\draw[qedge,-{Stealth[length=3pt]}] (9.22826,0.75499) -- (12.00000,0.00000);
\draw[qedge,-{Stealth[length=3pt]}] (10.01414,2.60858) -- (12.00000,0.00000);
\draw[qedge,-{Stealth[length=3pt]}] (2.69330,2.96621) -- (4.04994,0.42777);
\draw[qedge,-{Stealth[length=3pt]}] (6.24482,1.48372) -- (7.28809,3.25717);
\draw[qedge,-{Stealth[length=3pt]}] (5.23909,3.48455) -- (2.69330,2.96621);
\draw[qedge,-{Stealth[length=3pt]}] (7.28809,3.25717) -- (6.18899,7.50100);
\draw[qedge,-{Stealth[length=3pt]}] (6.18899,7.50100) -- (6.24482,1.48372);
\fill (0.00000,0.00000) circle[radius=1.1pt];
\draw[line width=.7pt,fill=white] (0.00000,0.00000) circle[radius=2.4pt];
\fill (12.00000,0.00000) circle[radius=1.1pt];
\draw[line width=.7pt,fill=white] (12.00000,0.00000) circle[radius=2.4pt];
\fill (6.00000,10.39200) circle[radius=1.1pt];
\draw[line width=.7pt,fill=white] (6.00000,10.39200) circle[radius=2.4pt];
\fill (4.04994,0.42777) circle[radius=1.1pt];
\fill (1.32113,1.43991) circle[radius=1.1pt];
\fill (10.01414,2.60858) circle[radius=1.1pt];
\fill (9.22826,0.75499) circle[radius=1.1pt];
\fill (2.69330,2.96621) circle[radius=1.1pt];
\fill (8.02049,5.26608) circle[radius=1.1pt];
\fill (7.28809,3.25717) circle[radius=1.1pt];
\fill (4.51516,5.61611) circle[radius=1.1pt];
\fill (6.18899,7.50100) circle[radius=1.1pt];
\fill (5.23909,3.48455) circle[radius=1.1pt];
\fill (6.24482,1.48372) circle[radius=1.1pt];
\node[anchor=north] at (6,-.25) {$t_0$ is the outer black face};
\end{tikzpicture}\\[1mm]$J(r,s)=3$\\[4mm]\begin{tikzpicture}[x=0.72cm,y=0.72cm]
\path[use as bounding box] (-.6,-.7) rectangle (12.6,11);
\fill[black!9] (0.00000,0.00000) -- (1.32113,1.43991) -- (4.04994,0.42777) -- cycle;
\fill[black!9] (6.00000,10.39200) -- (2.69330,2.96621) -- (1.32113,1.43991) -- cycle;
\fill[black!9] (6.00000,10.39200) -- (6.18899,7.50100) -- (4.51516,5.61611) -- cycle;
\fill[black!9] (6.00000,10.39200) -- (10.01414,2.60858) -- (8.02049,5.26608) -- cycle;
\fill[black!9] (12.00000,0.00000) -- (4.04994,0.42777) -- (9.22826,0.75499) -- cycle;
\fill[black!9] (12.00000,0.00000) -- (7.28809,3.25717) -- (10.01414,2.60858) -- cycle;
\fill[black!9] (4.04994,0.42777) -- (2.69330,2.96621) -- (6.24482,1.48372) -- cycle;
\fill[black!9] (9.22826,0.75499) -- (6.24482,1.48372) -- (7.28809,3.25717) -- cycle;
\fill[black!9] (2.69330,2.96621) -- (4.51516,5.61611) -- (5.23909,3.48455) -- cycle;
\fill[black!9] (8.02049,5.26608) -- (7.28809,3.25717) -- (6.18899,7.50100) -- cycle;
\fill[black!9] (6.18899,7.50100) -- (6.24482,1.48372) -- (5.23909,3.48455) -- cycle;
\draw[base] (0.00000,0.00000) -- (6.00000,10.39200);
\draw[base] (0.00000,0.00000) -- (12.00000,0.00000);
\draw[base] (0.00000,0.00000) -- (4.04994,0.42777);
\draw[base] (0.00000,0.00000) -- (1.32113,1.43991);
\draw[base] (6.00000,10.39200) -- (12.00000,0.00000);
\draw[base] (6.00000,10.39200) -- (1.32113,1.43991);
\draw[base] (6.00000,10.39200) -- (10.01414,2.60858);
\draw[base] (6.00000,10.39200) -- (2.69330,2.96621);
\draw[base] (6.00000,10.39200) -- (8.02049,5.26608);
\draw[base] (6.00000,10.39200) -- (4.51516,5.61611);
\draw[base] (6.00000,10.39200) -- (6.18899,7.50100);
\draw[base] (12.00000,0.00000) -- (4.04994,0.42777);
\draw[base] (12.00000,0.00000) -- (10.01414,2.60858);
\draw[base] (12.00000,0.00000) -- (9.22826,0.75499);
\draw[base] (12.00000,0.00000) -- (7.28809,3.25717);
\draw[base] (4.04994,0.42777) -- (1.32113,1.43991);
\draw[base] (4.04994,0.42777) -- (9.22826,0.75499);
\draw[base] (4.04994,0.42777) -- (2.69330,2.96621);
\draw[base] (4.04994,0.42777) -- (6.24482,1.48372);
\draw[base] (1.32113,1.43991) -- (2.69330,2.96621);
\draw[base] (10.01414,2.60858) -- (8.02049,5.26608);
\draw[base] (10.01414,2.60858) -- (7.28809,3.25717);
\draw[base] (9.22826,0.75499) -- (7.28809,3.25717);
\draw[base] (9.22826,0.75499) -- (6.24482,1.48372);
\draw[base] (2.69330,2.96621) -- (4.51516,5.61611);
\draw[base] (2.69330,2.96621) -- (5.23909,3.48455);
\draw[base] (2.69330,2.96621) -- (6.24482,1.48372);
\draw[base] (8.02049,5.26608) -- (7.28809,3.25717);
\draw[base] (8.02049,5.26608) -- (6.18899,7.50100);
\draw[base] (7.28809,3.25717) -- (6.18899,7.50100);
\draw[base] (7.28809,3.25717) -- (6.24482,1.48372);
\draw[base] (4.51516,5.61611) -- (6.18899,7.50100);
\draw[base] (4.51516,5.61611) -- (5.23909,3.48455);
\draw[base] (6.18899,7.50100) -- (5.23909,3.48455);
\draw[base] (6.18899,7.50100) -- (6.24482,1.48372);
\draw[base] (5.23909,3.48455) -- (6.24482,1.48372);
\node[markS] at (1.43374,1.24375) {};
\node[markR] at (3.50764,0.47452) {};
\node[markS] at (2.84806,3.43817) {};
\node[markR] at (1.80521,2.27818) {};
\node[markS] at (6.03996,7.58149) {};
\node[markR] at (4.76786,6.14897) {};
\node[markS] at (9.53351,3.44385) {};
\node[markR] at (8.01834,5.46356) {};
\node[markS] at (5.10021,0.41973) {};
\node[markR] at (9.03574,0.66841) {};
\node[markS] at (7.88313,2.94471) {};
\node[markR] at (9.95492,2.45178) {};
\node[markS] at (5.78511,1.51785) {};
\node[markR] at (3.08595,2.64453) {};
\node[markS] at (8.83437,1.01346) {};
\node[markR] at (7.35984,2.91512) {};
\node[markS] at (4.42733,5.23359) {};
\node[markR] at (4.97751,3.61361) {};
\node[markS] at (7.81538,5.28416) {};
\node[markR] at (6.42343,6.98270) {};
\node[markS] at (5.39554,3.64580) {};
\node[markR] at (6.15990,2.12517) {};
\draw[qedge,-{Stealth[length=3pt]}] (4.04994,0.42777) -- (0.00000,0.00000);
\draw[qedge,-{Stealth[length=3pt]}] (1.32113,1.43991) -- (6.00000,10.39200);
\draw[qedge,-{Stealth[length=3pt]}] (4.51516,5.61611) -- (6.00000,10.39200);
\draw[qedge,-{Stealth[length=3pt]}] (8.02049,5.26608) -- (6.00000,10.39200);
\draw[qedge,-{Stealth[length=3pt]}] (9.22826,0.75499) -- (12.00000,0.00000);
\draw[qedge,-{Stealth[length=3pt]}] (10.01414,2.60858) -- (12.00000,0.00000);
\draw[qedge,-{Stealth[length=3pt]}] (2.69330,2.96621) -- (4.04994,0.42777);
\draw[qedge,-{Stealth[length=3pt]}] (7.28809,3.25717) -- (6.24482,1.48372);
\draw[qedge,-{Stealth[length=3pt]}] (5.23909,3.48455) -- (2.69330,2.96621);
\draw[qedge,-{Stealth[length=3pt]}] (6.18899,7.50100) -- (7.28809,3.25717);
\draw[qedge,-{Stealth[length=3pt]}] (6.24482,1.48372) -- (6.18899,7.50100);
\fill (0.00000,0.00000) circle[radius=1.1pt];
\draw[line width=.7pt,fill=white] (0.00000,0.00000) circle[radius=2.4pt];
\fill (12.00000,0.00000) circle[radius=1.1pt];
\draw[line width=.7pt,fill=white] (12.00000,0.00000) circle[radius=2.4pt];
\fill (6.00000,10.39200) circle[radius=1.1pt];
\draw[line width=.7pt,fill=white] (6.00000,10.39200) circle[radius=2.4pt];
\fill (4.04994,0.42777) circle[radius=1.1pt];
\fill (1.32113,1.43991) circle[radius=1.1pt];
\fill (10.01414,2.60858) circle[radius=1.1pt];
\fill (9.22826,0.75499) circle[radius=1.1pt];
\fill (2.69330,2.96621) circle[radius=1.1pt];
\fill (8.02049,5.26608) circle[radius=1.1pt];
\fill (7.28809,3.25717) circle[radius=1.1pt];
\fill (4.51516,5.61611) circle[radius=1.1pt];
\fill (6.18899,7.50100) circle[radius=1.1pt];
\fill (5.23909,3.48455) circle[radius=1.1pt];
\fill (6.24482,1.48372) circle[radius=1.1pt];
\node[anchor=north] at (6,-.25) {$t_0$ is the outer black face};
\end{tikzpicture}\\[1mm]$J(\widetilde r,s)=1$
\end{tabular}
\caption{Reverse the dashed cycle by moving only the square marks. Blue disks, hence $s$ and its real weight, stay fixed. The phase changes from $i^3$ to $i$, so these two terms cancel. Applying the move twice restores the first panel.}
\end{figure}
\clearpage
\section*{9. spokes}
\begin{figure}[h!]\centering
\begin{tabular}{@{}c@{}}
\begin{tikzpicture}[x=0.72cm,y=0.72cm]
\path[use as bounding box] (-.6,-.7) rectangle (12.6,11);
\fill[black!9] (0.00000,0.00000) -- (1.32113,1.43991) -- (4.04994,0.42777) -- cycle;
\fill[black!9] (6.00000,10.39200) -- (2.69330,2.96621) -- (1.32113,1.43991) -- cycle;
\fill[black!9] (6.00000,10.39200) -- (6.18899,7.50100) -- (4.51516,5.61611) -- cycle;
\fill[black!9] (6.00000,10.39200) -- (10.01414,2.60858) -- (8.02049,5.26608) -- cycle;
\fill[black!9] (12.00000,0.00000) -- (4.04994,0.42777) -- (9.22826,0.75499) -- cycle;
\fill[black!9] (12.00000,0.00000) -- (7.28809,3.25717) -- (10.01414,2.60858) -- cycle;
\fill[black!9] (4.04994,0.42777) -- (2.69330,2.96621) -- (6.24482,1.48372) -- cycle;
\fill[black!9] (9.22826,0.75499) -- (6.24482,1.48372) -- (7.28809,3.25717) -- cycle;
\fill[black!9] (2.69330,2.96621) -- (4.51516,5.61611) -- (5.23909,3.48455) -- cycle;
\fill[black!9] (8.02049,5.26608) -- (7.28809,3.25717) -- (6.18899,7.50100) -- cycle;
\fill[black!9] (6.18899,7.50100) -- (6.24482,1.48372) -- (5.23909,3.48455) -- cycle;
\draw[base] (0.00000,0.00000) -- (6.00000,10.39200);
\draw[base] (0.00000,0.00000) -- (12.00000,0.00000);
\draw[base] (0.00000,0.00000) -- (4.04994,0.42777);
\draw[base] (0.00000,0.00000) -- (1.32113,1.43991);
\draw[base] (6.00000,10.39200) -- (12.00000,0.00000);
\draw[base] (6.00000,10.39200) -- (1.32113,1.43991);
\draw[base] (6.00000,10.39200) -- (10.01414,2.60858);
\draw[base] (6.00000,10.39200) -- (2.69330,2.96621);
\draw[base] (6.00000,10.39200) -- (8.02049,5.26608);
\draw[base] (6.00000,10.39200) -- (4.51516,5.61611);
\draw[base] (6.00000,10.39200) -- (6.18899,7.50100);
\draw[base] (12.00000,0.00000) -- (4.04994,0.42777);
\draw[base] (12.00000,0.00000) -- (10.01414,2.60858);
\draw[base] (12.00000,0.00000) -- (9.22826,0.75499);
\draw[base] (12.00000,0.00000) -- (7.28809,3.25717);
\draw[base] (4.04994,0.42777) -- (1.32113,1.43991);
\draw[base] (4.04994,0.42777) -- (9.22826,0.75499);
\draw[base] (4.04994,0.42777) -- (2.69330,2.96621);
\draw[base] (4.04994,0.42777) -- (6.24482,1.48372);
\draw[base] (1.32113,1.43991) -- (2.69330,2.96621);
\draw[base] (10.01414,2.60858) -- (8.02049,5.26608);
\draw[base] (10.01414,2.60858) -- (7.28809,3.25717);
\draw[base] (9.22826,0.75499) -- (7.28809,3.25717);
\draw[base] (9.22826,0.75499) -- (6.24482,1.48372);
\draw[base] (2.69330,2.96621) -- (4.51516,5.61611);
\draw[base] (2.69330,2.96621) -- (5.23909,3.48455);
\draw[base] (2.69330,2.96621) -- (6.24482,1.48372);
\draw[base] (8.02049,5.26608) -- (7.28809,3.25717);
\draw[base] (8.02049,5.26608) -- (6.18899,7.50100);
\draw[base] (7.28809,3.25717) -- (6.18899,7.50100);
\draw[base] (7.28809,3.25717) -- (6.24482,1.48372);
\draw[base] (4.51516,5.61611) -- (6.18899,7.50100);
\draw[base] (4.51516,5.61611) -- (5.23909,3.48455);
\draw[base] (6.18899,7.50100) -- (5.23909,3.48455);
\draw[base] (6.18899,7.50100) -- (6.24482,1.48372);
\draw[base] (5.23909,3.48455) -- (6.24482,1.48372);
\draw[dedge] (4.04994,0.42777) -- (1.32113,1.43991);
\draw[dedge] (1.32113,1.43991) -- (2.69330,2.96621);
\draw[dedge] (4.51516,5.61611) -- (6.18899,7.50100);
\draw[dedge] (8.02049,5.26608) -- (10.01414,2.60858);
\draw[dedge] (9.22826,0.75499) -- (4.04994,0.42777);
\draw[dedge] (10.01414,2.60858) -- (7.28809,3.25717);
\draw[dedge] (2.69330,2.96621) -- (6.24482,1.48372);
\draw[dedge] (6.24482,1.48372) -- (9.22826,0.75499);
\draw[dedge] (5.23909,3.48455) -- (4.51516,5.61611);
\draw[dedge] (7.28809,3.25717) -- (8.02049,5.26608);
\draw[dedge] (6.18899,7.50100) -- (5.23909,3.48455);
\fill (0.00000,0.00000) circle[radius=1.1pt];
\draw[line width=.7pt,fill=white] (0.00000,0.00000) circle[radius=2.4pt];
\fill (12.00000,0.00000) circle[radius=1.1pt];
\draw[line width=.7pt,fill=white] (12.00000,0.00000) circle[radius=2.4pt];
\fill (6.00000,10.39200) circle[radius=1.1pt];
\draw[line width=.7pt,fill=white] (6.00000,10.39200) circle[radius=2.4pt];
\fill (4.04994,0.42777) circle[radius=1.1pt];
\fill (1.32113,1.43991) circle[radius=1.1pt];
\fill (10.01414,2.60858) circle[radius=1.1pt];
\fill (9.22826,0.75499) circle[radius=1.1pt];
\fill (2.69330,2.96621) circle[radius=1.1pt];
\fill (8.02049,5.26608) circle[radius=1.1pt];
\fill (7.28809,3.25717) circle[radius=1.1pt];
\fill (4.51516,5.61611) circle[radius=1.1pt];
\fill (6.18899,7.50100) circle[radius=1.1pt];
\fill (5.23909,3.48455) circle[radius=1.1pt];
\fill (6.24482,1.48372) circle[radius=1.1pt];
\node[anchor=north] at (6,-.25) {$t_0$ is the outer black face};
\end{tikzpicture}\\[1mm]Edges of $D$\\[4mm]\begin{tikzpicture}[x=0.72cm,y=0.72cm]
\path[use as bounding box] (-.6,-.7) rectangle (12.6,11);
\fill[black!9] (0.00000,0.00000) -- (1.32113,1.43991) -- (4.04994,0.42777) -- cycle;
\fill[black!9] (6.00000,10.39200) -- (2.69330,2.96621) -- (1.32113,1.43991) -- cycle;
\fill[black!9] (6.00000,10.39200) -- (6.18899,7.50100) -- (4.51516,5.61611) -- cycle;
\fill[black!9] (6.00000,10.39200) -- (10.01414,2.60858) -- (8.02049,5.26608) -- cycle;
\fill[black!9] (12.00000,0.00000) -- (4.04994,0.42777) -- (9.22826,0.75499) -- cycle;
\fill[black!9] (12.00000,0.00000) -- (7.28809,3.25717) -- (10.01414,2.60858) -- cycle;
\fill[black!9] (4.04994,0.42777) -- (2.69330,2.96621) -- (6.24482,1.48372) -- cycle;
\fill[black!9] (9.22826,0.75499) -- (6.24482,1.48372) -- (7.28809,3.25717) -- cycle;
\fill[black!9] (2.69330,2.96621) -- (4.51516,5.61611) -- (5.23909,3.48455) -- cycle;
\fill[black!9] (8.02049,5.26608) -- (7.28809,3.25717) -- (6.18899,7.50100) -- cycle;
\fill[black!9] (6.18899,7.50100) -- (6.24482,1.48372) -- (5.23909,3.48455) -- cycle;
\draw[base] (0.00000,0.00000) -- (6.00000,10.39200);
\draw[base] (0.00000,0.00000) -- (12.00000,0.00000);
\draw[base] (0.00000,0.00000) -- (4.04994,0.42777);
\draw[base] (0.00000,0.00000) -- (1.32113,1.43991);
\draw[base] (6.00000,10.39200) -- (12.00000,0.00000);
\draw[base] (6.00000,10.39200) -- (1.32113,1.43991);
\draw[base] (6.00000,10.39200) -- (10.01414,2.60858);
\draw[base] (6.00000,10.39200) -- (2.69330,2.96621);
\draw[base] (6.00000,10.39200) -- (8.02049,5.26608);
\draw[base] (6.00000,10.39200) -- (4.51516,5.61611);
\draw[base] (6.00000,10.39200) -- (6.18899,7.50100);
\draw[base] (12.00000,0.00000) -- (4.04994,0.42777);
\draw[base] (12.00000,0.00000) -- (10.01414,2.60858);
\draw[base] (12.00000,0.00000) -- (9.22826,0.75499);
\draw[base] (12.00000,0.00000) -- (7.28809,3.25717);
\draw[base] (4.04994,0.42777) -- (1.32113,1.43991);
\draw[base] (4.04994,0.42777) -- (9.22826,0.75499);
\draw[base] (4.04994,0.42777) -- (2.69330,2.96621);
\draw[base] (4.04994,0.42777) -- (6.24482,1.48372);
\draw[base] (1.32113,1.43991) -- (2.69330,2.96621);
\draw[base] (10.01414,2.60858) -- (8.02049,5.26608);
\draw[base] (10.01414,2.60858) -- (7.28809,3.25717);
\draw[base] (9.22826,0.75499) -- (7.28809,3.25717);
\draw[base] (9.22826,0.75499) -- (6.24482,1.48372);
\draw[base] (2.69330,2.96621) -- (4.51516,5.61611);
\draw[base] (2.69330,2.96621) -- (5.23909,3.48455);
\draw[base] (2.69330,2.96621) -- (6.24482,1.48372);
\draw[base] (8.02049,5.26608) -- (7.28809,3.25717);
\draw[base] (8.02049,5.26608) -- (6.18899,7.50100);
\draw[base] (7.28809,3.25717) -- (6.18899,7.50100);
\draw[base] (7.28809,3.25717) -- (6.24482,1.48372);
\draw[base] (4.51516,5.61611) -- (6.18899,7.50100);
\draw[base] (4.51516,5.61611) -- (5.23909,3.48455);
\draw[base] (6.18899,7.50100) -- (5.23909,3.48455);
\draw[base] (6.18899,7.50100) -- (6.24482,1.48372);
\draw[base] (5.23909,3.48455) -- (6.24482,1.48372);
\draw[spoke] (1.79036,0.62256) -- (0.00000,0.00000);
\draw[spoke] (1.79036,0.62256) -- (1.32113,1.43991);
\draw[spoke] (1.79036,0.62256) -- (4.04994,0.42777);
\draw[spoke] (3.33814,4.93271) -- (6.00000,10.39200);
\draw[spoke] (3.33814,4.93271) -- (2.69330,2.96621);
\draw[spoke] (3.33814,4.93271) -- (1.32113,1.43991);
\draw[spoke] (5.56805,7.83637) -- (6.00000,10.39200);
\draw[spoke] (5.56805,7.83637) -- (6.18899,7.50100);
\draw[spoke] (5.56805,7.83637) -- (4.51516,5.61611);
\draw[spoke] (8.01154,6.08889) -- (6.00000,10.39200);
\draw[spoke] (8.01154,6.08889) -- (10.01414,2.60858);
\draw[spoke] (8.01154,6.08889) -- (8.02049,5.26608);
\draw[spoke] (8.42607,0.39425) -- (12.00000,0.00000);
\draw[spoke] (8.42607,0.39425) -- (4.04994,0.42777);
\draw[spoke] (8.42607,0.39425) -- (9.22826,0.75499);
\draw[spoke] (9.76741,1.95525) -- (12.00000,0.00000);
\draw[spoke] (9.76741,1.95525) -- (7.28809,3.25717);
\draw[spoke] (9.76741,1.95525) -- (10.01414,2.60858);
\draw[spoke] (4.32936,1.62590) -- (4.04994,0.42777);
\draw[spoke] (4.32936,1.62590) -- (2.69330,2.96621);
\draw[spoke] (4.32936,1.62590) -- (6.24482,1.48372);
\draw[spoke] (7.58706,1.83196) -- (9.22826,0.75499);
\draw[spoke] (7.58706,1.83196) -- (6.24482,1.48372);
\draw[spoke] (7.58706,1.83196) -- (7.28809,3.25717);
\draw[spoke] (4.14918,4.02229) -- (2.69330,2.96621);
\draw[spoke] (4.14918,4.02229) -- (4.51516,5.61611);
\draw[spoke] (4.14918,4.02229) -- (5.23909,3.48455);
\draw[spoke] (7.16586,5.34142) -- (8.02049,5.26608);
\draw[spoke] (7.16586,5.34142) -- (7.28809,3.25717);
\draw[spoke] (7.16586,5.34142) -- (6.18899,7.50100);
\draw[spoke] (5.89097,4.15642) -- (6.18899,7.50100);
\draw[spoke] (5.89097,4.15642) -- (6.24482,1.48372);
\draw[spoke] (5.89097,4.15642) -- (5.23909,3.48455);
\draw[dedge] (4.04994,0.42777) -- (1.79036,0.62256);
\draw[dedge] (1.79036,0.62256) -- (1.32113,1.43991);
\draw[dedge] (1.32113,1.43991) -- (3.33814,4.93271);
\draw[dedge] (3.33814,4.93271) -- (2.69330,2.96621);
\draw[dedge] (4.51516,5.61611) -- (5.56805,7.83637);
\draw[dedge] (5.56805,7.83637) -- (6.18899,7.50100);
\draw[dedge] (8.02049,5.26608) -- (8.01154,6.08889);
\draw[dedge] (8.01154,6.08889) -- (10.01414,2.60858);
\draw[dedge] (9.22826,0.75499) -- (8.42607,0.39425);
\draw[dedge] (8.42607,0.39425) -- (4.04994,0.42777);
\draw[dedge] (10.01414,2.60858) -- (9.76741,1.95525);
\draw[dedge] (9.76741,1.95525) -- (7.28809,3.25717);
\draw[dedge] (2.69330,2.96621) -- (4.32936,1.62590);
\draw[dedge] (4.32936,1.62590) -- (6.24482,1.48372);
\draw[dedge] (6.24482,1.48372) -- (7.58706,1.83196);
\draw[dedge] (7.58706,1.83196) -- (9.22826,0.75499);
\draw[dedge] (5.23909,3.48455) -- (4.14918,4.02229);
\draw[dedge] (4.14918,4.02229) -- (4.51516,5.61611);
\draw[dedge] (7.28809,3.25717) -- (7.16586,5.34142);
\draw[dedge] (7.16586,5.34142) -- (8.02049,5.26608);
\draw[dedge] (6.18899,7.50100) -- (5.89097,4.15642);
\draw[dedge] (5.89097,4.15642) -- (5.23909,3.48455);
\fill (0.00000,0.00000) circle[radius=1.1pt];
\draw[line width=.7pt,fill=white] (0.00000,0.00000) circle[radius=2.4pt];
\fill (12.00000,0.00000) circle[radius=1.1pt];
\draw[line width=.7pt,fill=white] (12.00000,0.00000) circle[radius=2.4pt];
\fill (6.00000,10.39200) circle[radius=1.1pt];
\draw[line width=.7pt,fill=white] (6.00000,10.39200) circle[radius=2.4pt];
\fill (4.04994,0.42777) circle[radius=1.1pt];
\fill (1.32113,1.43991) circle[radius=1.1pt];
\fill (10.01414,2.60858) circle[radius=1.1pt];
\fill (9.22826,0.75499) circle[radius=1.1pt];
\fill (2.69330,2.96621) circle[radius=1.1pt];
\fill (8.02049,5.26608) circle[radius=1.1pt];
\fill (7.28809,3.25717) circle[radius=1.1pt];
\fill (4.51516,5.61611) circle[radius=1.1pt];
\fill (6.18899,7.50100) circle[radius=1.1pt];
\fill (5.23909,3.48455) circle[radius=1.1pt];
\fill (6.24482,1.48372) circle[radius=1.1pt];
\node[anchor=north] at (6,-.25) {$t_0$ is the outer black face};
\end{tikzpicture}\\[1mm]Two-spoke detours
\end{tabular}
\caption{Each directed edge moves through the centre of its own black triangle. Detours have disjoint interiors and exclude the third corner. With $a(t,v)=\tfrac12\det(c_t,v)$, counterclockwise circulation is the positive area of the spoke polygon. The checked circulations are 5.641, 5.289, 11.033.}
\end{figure}
\clearpage
\section*{10. white tree}
\begin{figure}[h!]\centering
\begin{tabular}{@{}c@{}}
\begin{tikzpicture}[x=0.72cm,y=0.72cm]
\path[use as bounding box] (-1.5,-2.8) rectangle (14.5,12);
\fill[black!9] (0.00000,0.00000) -- (1.32113,1.43991) -- (4.04994,0.42777) -- cycle;
\fill[black!9] (6.00000,10.39200) -- (2.69330,2.96621) -- (1.32113,1.43991) -- cycle;
\fill[black!9] (6.00000,10.39200) -- (6.18899,7.50100) -- (4.51516,5.61611) -- cycle;
\fill[black!9] (6.00000,10.39200) -- (10.01414,2.60858) -- (8.02049,5.26608) -- cycle;
\fill[black!9] (12.00000,0.00000) -- (4.04994,0.42777) -- (9.22826,0.75499) -- cycle;
\fill[black!9] (12.00000,0.00000) -- (7.28809,3.25717) -- (10.01414,2.60858) -- cycle;
\fill[black!9] (4.04994,0.42777) -- (2.69330,2.96621) -- (6.24482,1.48372) -- cycle;
\fill[black!9] (9.22826,0.75499) -- (6.24482,1.48372) -- (7.28809,3.25717) -- cycle;
\fill[black!9] (2.69330,2.96621) -- (4.51516,5.61611) -- (5.23909,3.48455) -- cycle;
\fill[black!9] (8.02049,5.26608) -- (7.28809,3.25717) -- (6.18899,7.50100) -- cycle;
\fill[black!9] (6.18899,7.50100) -- (6.24482,1.48372) -- (5.23909,3.48455) -- cycle;
\draw[base] (0.00000,0.00000) -- (6.00000,10.39200);
\draw[base] (0.00000,0.00000) -- (12.00000,0.00000);
\draw[base] (0.00000,0.00000) -- (4.04994,0.42777);
\draw[base] (0.00000,0.00000) -- (1.32113,1.43991);
\draw[base] (6.00000,10.39200) -- (12.00000,0.00000);
\draw[base] (6.00000,10.39200) -- (1.32113,1.43991);
\draw[base] (6.00000,10.39200) -- (10.01414,2.60858);
\draw[base] (6.00000,10.39200) -- (2.69330,2.96621);
\draw[base] (6.00000,10.39200) -- (8.02049,5.26608);
\draw[base] (6.00000,10.39200) -- (4.51516,5.61611);
\draw[base] (6.00000,10.39200) -- (6.18899,7.50100);
\draw[base] (12.00000,0.00000) -- (4.04994,0.42777);
\draw[base] (12.00000,0.00000) -- (10.01414,2.60858);
\draw[base] (12.00000,0.00000) -- (9.22826,0.75499);
\draw[base] (12.00000,0.00000) -- (7.28809,3.25717);
\draw[base] (4.04994,0.42777) -- (1.32113,1.43991);
\draw[base] (4.04994,0.42777) -- (9.22826,0.75499);
\draw[base] (4.04994,0.42777) -- (2.69330,2.96621);
\draw[base] (4.04994,0.42777) -- (6.24482,1.48372);
\draw[base] (1.32113,1.43991) -- (2.69330,2.96621);
\draw[base] (10.01414,2.60858) -- (8.02049,5.26608);
\draw[base] (10.01414,2.60858) -- (7.28809,3.25717);
\draw[base] (9.22826,0.75499) -- (7.28809,3.25717);
\draw[base] (9.22826,0.75499) -- (6.24482,1.48372);
\draw[base] (2.69330,2.96621) -- (4.51516,5.61611);
\draw[base] (2.69330,2.96621) -- (5.23909,3.48455);
\draw[base] (2.69330,2.96621) -- (6.24482,1.48372);
\draw[base] (8.02049,5.26608) -- (7.28809,3.25717);
\draw[base] (8.02049,5.26608) -- (6.18899,7.50100);
\draw[base] (7.28809,3.25717) -- (6.18899,7.50100);
\draw[base] (7.28809,3.25717) -- (6.24482,1.48372);
\draw[base] (4.51516,5.61611) -- (6.18899,7.50100);
\draw[base] (4.51516,5.61611) -- (5.23909,3.48455);
\draw[base] (6.18899,7.50100) -- (5.23909,3.48455);
\draw[base] (6.18899,7.50100) -- (6.24482,1.48372);
\draw[base] (5.23909,3.48455) -- (6.24482,1.48372);
\draw[pedge] (1.79036,0.62256) -- (0.00000,0.00000);
\draw[pedge] (1.79036,0.62256) -- (1.32113,1.43991);
\draw[spoke,dashed] (1.79036,0.62256) -- (4.04994,0.42777);
\fill (1.79036,0.62256) circle[radius=1.5pt];
\draw[pedge] (3.33814,4.93271) -- (6.00000,10.39200);
\draw[pedge] (3.33814,4.93271) -- (2.69330,2.96621);
\draw[spoke,dashed] (3.33814,4.93271) -- (1.32113,1.43991);
\fill (3.33814,4.93271) circle[radius=1.5pt];
\draw[pedge] (5.56805,7.83637) -- (6.00000,10.39200);
\draw[pedge] (5.56805,7.83637) -- (6.18899,7.50100);
\draw[spoke,dashed] (5.56805,7.83637) -- (4.51516,5.61611);
\fill (5.56805,7.83637) circle[radius=1.5pt];
\draw[pedge] (8.01154,6.08889) -- (6.00000,10.39200);
\draw[pedge] (8.01154,6.08889) -- (10.01414,2.60858);
\draw[spoke,dashed] (8.01154,6.08889) -- (8.02049,5.26608);
\fill (8.01154,6.08889) circle[radius=1.5pt];
\draw[pedge] (8.42607,0.39425) -- (12.00000,0.00000);
\draw[pedge] (8.42607,0.39425) -- (4.04994,0.42777);
\draw[spoke,dashed] (8.42607,0.39425) -- (9.22826,0.75499);
\fill (8.42607,0.39425) circle[radius=1.5pt];
\draw[pedge] (9.76741,1.95525) -- (12.00000,0.00000);
\draw[pedge] (9.76741,1.95525) -- (7.28809,3.25717);
\draw[spoke,dashed] (9.76741,1.95525) -- (10.01414,2.60858);
\fill (9.76741,1.95525) circle[radius=1.5pt];
\draw[pedge] (4.32936,1.62590) -- (4.04994,0.42777);
\draw[spoke,dashed] (4.32936,1.62590) -- (2.69330,2.96621);
\draw[pedge] (4.32936,1.62590) -- (6.24482,1.48372);
\fill (4.32936,1.62590) circle[radius=1.5pt];
\draw[pedge] (7.58706,1.83196) -- (9.22826,0.75499);
\draw[spoke,dashed] (7.58706,1.83196) -- (6.24482,1.48372);
\draw[pedge] (7.58706,1.83196) -- (7.28809,3.25717);
\fill (7.58706,1.83196) circle[radius=1.5pt];
\draw[pedge] (4.14918,4.02229) -- (2.69330,2.96621);
\draw[pedge] (4.14918,4.02229) -- (4.51516,5.61611);
\draw[spoke,dashed] (4.14918,4.02229) -- (5.23909,3.48455);
\fill (4.14918,4.02229) circle[radius=1.5pt];
\draw[pedge] (7.16586,5.34142) -- (8.02049,5.26608);
\draw[spoke,dashed] (7.16586,5.34142) -- (7.28809,3.25717);
\draw[pedge] (7.16586,5.34142) -- (6.18899,7.50100);
\fill (7.16586,5.34142) circle[radius=1.5pt];
\draw[spoke,dashed] (5.89097,4.15642) -- (6.18899,7.50100);
\draw[pedge] (5.89097,4.15642) -- (6.24482,1.48372);
\draw[pedge] (5.89097,4.15642) -- (5.23909,3.48455);
\fill (5.89097,4.15642) circle[radius=1.5pt];
\draw[pedge] (6.00000,-1.40000) -- (0.00000,0.00000);
\draw[pedge] (6.00000,-1.40000) -- (12.00000,0.00000);
\draw[pedge] (6,-1.4)--(14,-1.4)--(14,11.5)--(6,11.5)--(6,10.392);
\node[fill=white,inner sep=1pt] at (6,-1.4) {$c_0$};
\fill (0.00000,0.00000) circle[radius=1.1pt];
\draw[line width=.7pt,fill=white] (0.00000,0.00000) circle[radius=2.4pt];
\fill (12.00000,0.00000) circle[radius=1.1pt];
\draw[line width=.7pt,fill=white] (12.00000,0.00000) circle[radius=2.4pt];
\fill (6.00000,10.39200) circle[radius=1.1pt];
\draw[line width=.7pt,fill=white] (6.00000,10.39200) circle[radius=2.4pt];
\fill (4.04994,0.42777) circle[radius=1.1pt];
\fill (1.32113,1.43991) circle[radius=1.1pt];
\fill (10.01414,2.60858) circle[radius=1.1pt];
\fill (9.22826,0.75499) circle[radius=1.1pt];
\fill (2.69330,2.96621) circle[radius=1.1pt];
\fill (8.02049,5.26608) circle[radius=1.1pt];
\fill (7.28809,3.25717) circle[radius=1.1pt];
\fill (4.51516,5.61611) circle[radius=1.1pt];
\fill (6.18899,7.50100) circle[radius=1.1pt];
\fill (5.23909,3.48455) circle[radius=1.1pt];
\fill (6.24482,1.48372) circle[radius=1.1pt];
\node[anchor=north] at (6,-2.3) {$t_0$ is the outer black face};
\end{tikzpicture}\\[1mm]Retained incidence tree $H-s$\\[4mm]\begin{tikzpicture}[x=0.72cm,y=0.72cm]
\path[use as bounding box] (-1.5,-1.4) rectangle (13.5,11);
\fill[black!9] (0.00000,0.00000) -- (1.32113,1.43991) -- (4.04994,0.42777) -- cycle;
\fill[black!9] (6.00000,10.39200) -- (2.69330,2.96621) -- (1.32113,1.43991) -- cycle;
\fill[black!9] (6.00000,10.39200) -- (6.18899,7.50100) -- (4.51516,5.61611) -- cycle;
\fill[black!9] (6.00000,10.39200) -- (10.01414,2.60858) -- (8.02049,5.26608) -- cycle;
\fill[black!9] (12.00000,0.00000) -- (4.04994,0.42777) -- (9.22826,0.75499) -- cycle;
\fill[black!9] (12.00000,0.00000) -- (7.28809,3.25717) -- (10.01414,2.60858) -- cycle;
\fill[black!9] (4.04994,0.42777) -- (2.69330,2.96621) -- (6.24482,1.48372) -- cycle;
\fill[black!9] (9.22826,0.75499) -- (6.24482,1.48372) -- (7.28809,3.25717) -- cycle;
\fill[black!9] (2.69330,2.96621) -- (4.51516,5.61611) -- (5.23909,3.48455) -- cycle;
\fill[black!9] (8.02049,5.26608) -- (7.28809,3.25717) -- (6.18899,7.50100) -- cycle;
\fill[black!9] (6.18899,7.50100) -- (6.24482,1.48372) -- (5.23909,3.48455) -- cycle;
\draw[base] (0.00000,0.00000) -- (6.00000,10.39200);
\draw[base] (0.00000,0.00000) -- (12.00000,0.00000);
\draw[base] (0.00000,0.00000) -- (4.04994,0.42777);
\draw[base] (0.00000,0.00000) -- (1.32113,1.43991);
\draw[base] (6.00000,10.39200) -- (12.00000,0.00000);
\draw[base] (6.00000,10.39200) -- (1.32113,1.43991);
\draw[base] (6.00000,10.39200) -- (10.01414,2.60858);
\draw[base] (6.00000,10.39200) -- (2.69330,2.96621);
\draw[base] (6.00000,10.39200) -- (8.02049,5.26608);
\draw[base] (6.00000,10.39200) -- (4.51516,5.61611);
\draw[base] (6.00000,10.39200) -- (6.18899,7.50100);
\draw[base] (12.00000,0.00000) -- (4.04994,0.42777);
\draw[base] (12.00000,0.00000) -- (10.01414,2.60858);
\draw[base] (12.00000,0.00000) -- (9.22826,0.75499);
\draw[base] (12.00000,0.00000) -- (7.28809,3.25717);
\draw[base] (4.04994,0.42777) -- (1.32113,1.43991);
\draw[base] (4.04994,0.42777) -- (9.22826,0.75499);
\draw[base] (4.04994,0.42777) -- (2.69330,2.96621);
\draw[base] (4.04994,0.42777) -- (6.24482,1.48372);
\draw[base] (1.32113,1.43991) -- (2.69330,2.96621);
\draw[base] (10.01414,2.60858) -- (8.02049,5.26608);
\draw[base] (10.01414,2.60858) -- (7.28809,3.25717);
\draw[base] (9.22826,0.75499) -- (7.28809,3.25717);
\draw[base] (9.22826,0.75499) -- (6.24482,1.48372);
\draw[base] (2.69330,2.96621) -- (4.51516,5.61611);
\draw[base] (2.69330,2.96621) -- (5.23909,3.48455);
\draw[base] (2.69330,2.96621) -- (6.24482,1.48372);
\draw[base] (8.02049,5.26608) -- (7.28809,3.25717);
\draw[base] (8.02049,5.26608) -- (6.18899,7.50100);
\draw[base] (7.28809,3.25717) -- (6.18899,7.50100);
\draw[base] (7.28809,3.25717) -- (6.24482,1.48372);
\draw[base] (4.51516,5.61611) -- (6.18899,7.50100);
\draw[base] (4.51516,5.61611) -- (5.23909,3.48455);
\draw[base] (6.18899,7.50100) -- (5.23909,3.48455);
\draw[base] (6.18899,7.50100) -- (6.24482,1.48372);
\draw[base] (5.23909,3.48455) -- (6.24482,1.48372);
\draw[wt] (5.34998,0.14259) -- (2.02497,0.21389) -- (2.92015,0.52517) -- (2.68554,0.93384) -- (2.68812,1.61130);
\draw[wt] (2.44038,3.94397) -- (3.66057,5.91596) -- (2.32964,3.18631) -- (2.00722,2.20306) -- (2.68812,1.61130);
\draw[wt] (4.40282,6.32477) -- (5.25758,8.00405) -- (5.04161,6.72624) -- (5.35208,6.55855) -- (5.31441,5.53389);
\draw[wt] (6.73649,7.71969) -- (7.01025,7.82904) -- (8.01602,5.67749) -- (9.01731,3.93733) -- (8.44090,3.71061);
\draw[wt] (9.50545,1.33739) -- (10.61413,0.37749) -- (8.82717,0.57462) -- (6.63910,0.59138) -- (6.50768,0.88883);
\draw[wt] (9.33805,4.33353) -- (11.00707,1.30429) -- (9.89077,2.28191) -- (8.65111,2.93287) -- (8.44090,3.71061);
\draw[wt] (2.68812,1.61130) -- (3.37162,1.69699) -- (3.51133,2.29605) -- (4.46906,2.22497) -- (4.72574,2.64483);
\draw[wt] (6.50768,0.88883) -- (7.73654,1.11936) -- (6.91594,1.65784) -- (6.76646,2.37045) -- (6.57397,4.08063);
\draw[wt] (4.72574,2.64483) -- (3.96619,3.22538) -- (4.69414,3.75342) -- (4.87713,4.55033) -- (5.31441,5.53389);
\draw[wt] (8.44090,3.71061) -- (7.65429,4.26163) -- (7.22697,4.29929) -- (6.73854,5.37908) -- (6.57397,4.08063);
\draw[wt] (6.57397,4.08063) -- (6.21690,4.49236) -- (6.03998,5.82871) -- (5.71404,5.49277) -- (5.31441,5.53389);
\node[draw=treegreen,circle,line width=1pt,inner sep=3pt] at (5.34998,0.14259) {};
\node[anchor=south,fill=white,inner sep=.3pt] at (5.54998,0.29259) {$y_0$};
\node[circle,draw=treegreen,fill=white,inner sep=1.5pt] at (5.34998,0.14259) {};
\node[circle,draw=treegreen,fill=white,inner sep=1.5pt] at (2.44038,3.94397) {};
\node[circle,draw=treegreen,fill=white,inner sep=1.5pt] at (9.33805,4.33353) {};
\node[circle,draw=treegreen,fill=white,inner sep=1.5pt] at (4.40282,6.32477) {};
\node[circle,draw=treegreen,fill=white,inner sep=1.5pt] at (6.73649,7.71969) {};
\node[circle,draw=treegreen,fill=white,inner sep=1.5pt] at (9.50545,1.33739) {};
\node[circle,draw=treegreen,fill=white,inner sep=1.5pt] at (2.68812,1.61130) {};
\node[circle,draw=treegreen,fill=white,inner sep=1.5pt] at (6.50768,0.88883) {};
\node[circle,draw=treegreen,fill=white,inner sep=1.5pt] at (8.44090,3.71061) {};
\node[circle,draw=treegreen,fill=white,inner sep=1.5pt] at (4.72574,2.64483) {};
\node[circle,draw=treegreen,fill=white,inner sep=1.5pt] at (6.57397,4.08063) {};
\node[circle,draw=treegreen,fill=white,inner sep=1.5pt] at (5.31441,5.53389) {};
\fill (0.00000,0.00000) circle[radius=1.1pt];
\draw[line width=.7pt,fill=white] (0.00000,0.00000) circle[radius=2.4pt];
\fill (12.00000,0.00000) circle[radius=1.1pt];
\draw[line width=.7pt,fill=white] (12.00000,0.00000) circle[radius=2.4pt];
\fill (6.00000,10.39200) circle[radius=1.1pt];
\draw[line width=.7pt,fill=white] (6.00000,10.39200) circle[radius=2.4pt];
\fill (4.04994,0.42777) circle[radius=1.1pt];
\fill (1.32113,1.43991) circle[radius=1.1pt];
\fill (10.01414,2.60858) circle[radius=1.1pt];
\fill (9.22826,0.75499) circle[radius=1.1pt];
\fill (2.69330,2.96621) circle[radius=1.1pt];
\fill (8.02049,5.26608) circle[radius=1.1pt];
\fill (7.28809,3.25717) circle[radius=1.1pt];
\fill (4.51516,5.61611) circle[radius=1.1pt];
\fill (6.18899,7.50100) circle[radius=1.1pt];
\fill (5.23909,3.48455) circle[radius=1.1pt];
\fill (6.24482,1.48372) circle[radius=1.1pt];
\node[anchor=north] at (6,-.25) {$t_0$ is the outer black face};
\end{tikzpicture}\\[1mm]White-face cotree, rooted at $y_0$
\end{tabular}
\caption{The dash-dot routes cross the deleted spokes $(t,s(t))$. They join all 12 white centres in one tree with 11 edges. Each route stays within the two white faces and their common black triangle. This is the cotree of the retained incidence tree.}
\end{figure}
\clearpage
\section*{11. selected Jordan}
\begin{figure}[h!]\centering
\begin{tikzpicture}[x=0.72cm,y=0.72cm]
\path[use as bounding box] (-1.5,-1.4) rectangle (13.5,11);
\fill[black!9] (0.00000,0.00000) -- (1.32113,1.43991) -- (4.04994,0.42777) -- cycle;
\fill[black!9] (6.00000,10.39200) -- (2.69330,2.96621) -- (1.32113,1.43991) -- cycle;
\fill[black!9] (6.00000,10.39200) -- (6.18899,7.50100) -- (4.51516,5.61611) -- cycle;
\fill[black!9] (6.00000,10.39200) -- (10.01414,2.60858) -- (8.02049,5.26608) -- cycle;
\fill[black!9] (12.00000,0.00000) -- (4.04994,0.42777) -- (9.22826,0.75499) -- cycle;
\fill[black!9] (12.00000,0.00000) -- (7.28809,3.25717) -- (10.01414,2.60858) -- cycle;
\fill[black!9] (4.04994,0.42777) -- (2.69330,2.96621) -- (6.24482,1.48372) -- cycle;
\fill[black!9] (9.22826,0.75499) -- (6.24482,1.48372) -- (7.28809,3.25717) -- cycle;
\fill[black!9] (2.69330,2.96621) -- (4.51516,5.61611) -- (5.23909,3.48455) -- cycle;
\fill[black!9] (8.02049,5.26608) -- (7.28809,3.25717) -- (6.18899,7.50100) -- cycle;
\fill[black!9] (6.18899,7.50100) -- (6.24482,1.48372) -- (5.23909,3.48455) -- cycle;
\draw[base] (0.00000,0.00000) -- (6.00000,10.39200);
\draw[base] (0.00000,0.00000) -- (12.00000,0.00000);
\draw[base] (0.00000,0.00000) -- (4.04994,0.42777);
\draw[base] (0.00000,0.00000) -- (1.32113,1.43991);
\draw[base] (6.00000,10.39200) -- (12.00000,0.00000);
\draw[base] (6.00000,10.39200) -- (1.32113,1.43991);
\draw[base] (6.00000,10.39200) -- (10.01414,2.60858);
\draw[base] (6.00000,10.39200) -- (2.69330,2.96621);
\draw[base] (6.00000,10.39200) -- (8.02049,5.26608);
\draw[base] (6.00000,10.39200) -- (4.51516,5.61611);
\draw[base] (6.00000,10.39200) -- (6.18899,7.50100);
\draw[base] (12.00000,0.00000) -- (4.04994,0.42777);
\draw[base] (12.00000,0.00000) -- (10.01414,2.60858);
\draw[base] (12.00000,0.00000) -- (9.22826,0.75499);
\draw[base] (12.00000,0.00000) -- (7.28809,3.25717);
\draw[base] (4.04994,0.42777) -- (1.32113,1.43991);
\draw[base] (4.04994,0.42777) -- (9.22826,0.75499);
\draw[base] (4.04994,0.42777) -- (2.69330,2.96621);
\draw[base] (4.04994,0.42777) -- (6.24482,1.48372);
\draw[base] (1.32113,1.43991) -- (2.69330,2.96621);
\draw[base] (10.01414,2.60858) -- (8.02049,5.26608);
\draw[base] (10.01414,2.60858) -- (7.28809,3.25717);
\draw[base] (9.22826,0.75499) -- (7.28809,3.25717);
\draw[base] (9.22826,0.75499) -- (6.24482,1.48372);
\draw[base] (2.69330,2.96621) -- (4.51516,5.61611);
\draw[base] (2.69330,2.96621) -- (5.23909,3.48455);
\draw[base] (2.69330,2.96621) -- (6.24482,1.48372);
\draw[base] (8.02049,5.26608) -- (7.28809,3.25717);
\draw[base] (8.02049,5.26608) -- (6.18899,7.50100);
\draw[base] (7.28809,3.25717) -- (6.18899,7.50100);
\draw[base] (7.28809,3.25717) -- (6.24482,1.48372);
\draw[base] (4.51516,5.61611) -- (6.18899,7.50100);
\draw[base] (4.51516,5.61611) -- (5.23909,3.48455);
\draw[base] (6.18899,7.50100) -- (5.23909,3.48455);
\draw[base] (6.18899,7.50100) -- (6.24482,1.48372);
\draw[base] (5.23909,3.48455) -- (6.24482,1.48372);
\draw[pedge] (0.00000,0.00000) -- (12.00000,0.00000);
\draw[pedge] (4.04994,0.42777) -- (1.32113,1.43991);
\draw[pedge] (6.00000,10.39200) -- (1.32113,1.43991);
\draw[pedge] (6.00000,10.39200) -- (4.51516,5.61611);
\draw[pedge] (6.00000,10.39200) -- (8.02049,5.26608);
\draw[pedge] (12.00000,0.00000) -- (9.22826,0.75499);
\draw[pedge] (12.00000,0.00000) -- (10.01414,2.60858);
\draw[pedge] (2.69330,2.96621) -- (6.24482,1.48372);
\draw[pedge] (9.22826,0.75499) -- (6.24482,1.48372);
\draw[pedge] (4.51516,5.61611) -- (5.23909,3.48455);
\draw[pedge] (8.02049,5.26608) -- (7.28809,3.25717);
\draw[pedge] (6.18899,7.50100) -- (6.24482,1.48372);
\draw[ham] (2.44038,3.94397) -- (3.00000,5.19600) -- (1.70000,5.69600) -- (-1.00000,-1.00000) -- (6.00000,-1.00000);
\draw[ham] (9.33805,4.33353) -- (9.00000,5.19600) -- (10.30000,5.69600) -- (13.00000,-1.00000) -- (6.00000,-1.00000);
\draw[ham] (5.34998,0.14259) -- (2.02497,0.21389) -- (1.79036,0.62256);
\draw[ham] (5.34998,0.14259) -- (8.02497,0.21389) -- (8.42607,0.39425);
\draw[ham] (1.79036,0.62256) -- (0.66057,0.71996) -- (2.44038,3.94397);
\draw[ham] (9.33805,4.33353) -- (8.00707,6.50029) -- (8.01154,6.08889);
\draw[ham] (3.33814,4.93271) -- (4.34665,6.67910) -- (4.40282,6.32477);
\draw[ham] (3.33814,4.93271) -- (2.00722,2.20306) -- (2.68812,1.61130);
\draw[ham] (4.40282,6.32477) -- (3.60423,4.29116) -- (4.14918,4.02229);
\draw[ham] (5.56805,7.83637) -- (6.09449,8.94650) -- (6.73649,7.71969);
\draw[ham] (5.56805,7.83637) -- (5.35208,6.55855) -- (5.31441,5.53389);
\draw[ham] (6.73649,7.71969) -- (7.10474,6.38354) -- (7.16586,5.34142);
\draw[ham] (8.01154,6.08889) -- (9.01731,3.93733) -- (8.44090,3.71061);
\draw[ham] (8.42607,0.39425) -- (6.63910,0.59138) -- (6.50768,0.88883);
\draw[ham] (9.76741,1.95525) -- (9.64404,1.62858) -- (9.50545,1.33739);
\draw[ham] (9.76741,1.95525) -- (8.65111,2.93287) -- (8.44090,3.71061);
\draw[ham] (9.50545,1.33739) -- (8.25818,2.00608) -- (7.58706,1.83196);
\draw[ham] (2.68812,1.61130) -- (3.37162,1.69699) -- (4.32936,1.62590);
\draw[ham] (6.50768,0.88883) -- (5.14738,0.95575) -- (4.32936,1.62590);
\draw[ham] (7.58706,1.83196) -- (6.76646,2.37045) -- (6.57397,4.08063);
\draw[ham] (4.14918,4.02229) -- (3.96619,3.22538) -- (4.72574,2.64483);
\draw[ham] (4.72574,2.64483) -- (5.74196,2.48414) -- (5.89097,4.15642);
\draw[ham] (7.16586,5.34142) -- (6.73854,5.37908) -- (6.57397,4.08063);
\draw[ham] (5.31441,5.53389) -- (5.71404,5.49277) -- (5.89097,4.15642);
\fill (0.00000,0.00000) circle[radius=1.1pt];
\draw[line width=.7pt,fill=white] (0.00000,0.00000) circle[radius=2.4pt];
\fill (12.00000,0.00000) circle[radius=1.1pt];
\draw[line width=.7pt,fill=white] (12.00000,0.00000) circle[radius=2.4pt];
\fill (6.00000,10.39200) circle[radius=1.1pt];
\draw[line width=.7pt,fill=white] (6.00000,10.39200) circle[radius=2.4pt];
\fill (4.04994,0.42777) circle[radius=1.1pt];
\fill (1.32113,1.43991) circle[radius=1.1pt];
\fill (10.01414,2.60858) circle[radius=1.1pt];
\fill (9.22826,0.75499) circle[radius=1.1pt];
\fill (2.69330,2.96621) circle[radius=1.1pt];
\fill (8.02049,5.26608) circle[radius=1.1pt];
\fill (7.28809,3.25717) circle[radius=1.1pt];
\fill (4.51516,5.61611) circle[radius=1.1pt];
\fill (6.18899,7.50100) circle[radius=1.1pt];
\fill (5.23909,3.48455) circle[radius=1.1pt];
\fill (6.24482,1.48372) circle[radius=1.1pt];
\node[anchor=north] at (6,-.25) {$t_0$ is the outer black face};
\end{tikzpicture}
\caption{The bold black edges are $P$, one side of every face. The blue Jordan curve crosses precisely the unselected edges. Its two regions separate the two components of $P$; the exterior arcs represent edges incident with the outer-face dual vertex.}
\end{figure}
\clearpage
\section*{12. recovered}
\begin{figure}[h!]\centering
\begin{tabular}{@{}c@{}}
\begin{tikzpicture}[x=0.72cm,y=0.72cm]
\path[use as bounding box] (-.6,-.7) rectangle (12.6,11);
\fill[black!9] (0.00000,0.00000) -- (1.32113,1.43991) -- (4.04994,0.42777) -- cycle;
\fill[black!9] (6.00000,10.39200) -- (2.69330,2.96621) -- (1.32113,1.43991) -- cycle;
\fill[black!9] (6.00000,10.39200) -- (6.18899,7.50100) -- (4.51516,5.61611) -- cycle;
\fill[black!9] (6.00000,10.39200) -- (10.01414,2.60858) -- (8.02049,5.26608) -- cycle;
\fill[black!9] (12.00000,0.00000) -- (4.04994,0.42777) -- (9.22826,0.75499) -- cycle;
\fill[black!9] (12.00000,0.00000) -- (7.28809,3.25717) -- (10.01414,2.60858) -- cycle;
\fill[black!9] (4.04994,0.42777) -- (2.69330,2.96621) -- (6.24482,1.48372) -- cycle;
\fill[black!9] (9.22826,0.75499) -- (6.24482,1.48372) -- (7.28809,3.25717) -- cycle;
\fill[black!9] (2.69330,2.96621) -- (4.51516,5.61611) -- (5.23909,3.48455) -- cycle;
\fill[black!9] (8.02049,5.26608) -- (7.28809,3.25717) -- (6.18899,7.50100) -- cycle;
\fill[black!9] (6.18899,7.50100) -- (6.24482,1.48372) -- (5.23909,3.48455) -- cycle;
\draw[base] (0.00000,0.00000) -- (6.00000,10.39200);
\draw[base] (0.00000,0.00000) -- (12.00000,0.00000);
\draw[base] (0.00000,0.00000) -- (4.04994,0.42777);
\draw[base] (0.00000,0.00000) -- (1.32113,1.43991);
\draw[base] (6.00000,10.39200) -- (12.00000,0.00000);
\draw[base] (6.00000,10.39200) -- (1.32113,1.43991);
\draw[base] (6.00000,10.39200) -- (10.01414,2.60858);
\draw[base] (6.00000,10.39200) -- (2.69330,2.96621);
\draw[base] (6.00000,10.39200) -- (8.02049,5.26608);
\draw[base] (6.00000,10.39200) -- (4.51516,5.61611);
\draw[base] (6.00000,10.39200) -- (6.18899,7.50100);
\draw[base] (12.00000,0.00000) -- (4.04994,0.42777);
\draw[base] (12.00000,0.00000) -- (10.01414,2.60858);
\draw[base] (12.00000,0.00000) -- (9.22826,0.75499);
\draw[base] (12.00000,0.00000) -- (7.28809,3.25717);
\draw[base] (4.04994,0.42777) -- (1.32113,1.43991);
\draw[base] (4.04994,0.42777) -- (9.22826,0.75499);
\draw[base] (4.04994,0.42777) -- (2.69330,2.96621);
\draw[base] (4.04994,0.42777) -- (6.24482,1.48372);
\draw[base] (1.32113,1.43991) -- (2.69330,2.96621);
\draw[base] (10.01414,2.60858) -- (8.02049,5.26608);
\draw[base] (10.01414,2.60858) -- (7.28809,3.25717);
\draw[base] (9.22826,0.75499) -- (7.28809,3.25717);
\draw[base] (9.22826,0.75499) -- (6.24482,1.48372);
\draw[base] (2.69330,2.96621) -- (4.51516,5.61611);
\draw[base] (2.69330,2.96621) -- (5.23909,3.48455);
\draw[base] (2.69330,2.96621) -- (6.24482,1.48372);
\draw[base] (8.02049,5.26608) -- (7.28809,3.25717);
\draw[base] (8.02049,5.26608) -- (6.18899,7.50100);
\draw[base] (7.28809,3.25717) -- (6.18899,7.50100);
\draw[base] (7.28809,3.25717) -- (6.24482,1.48372);
\draw[base] (4.51516,5.61611) -- (6.18899,7.50100);
\draw[base] (4.51516,5.61611) -- (5.23909,3.48455);
\draw[base] (6.18899,7.50100) -- (5.23909,3.48455);
\draw[base] (6.18899,7.50100) -- (6.24482,1.48372);
\draw[base] (5.23909,3.48455) -- (6.24482,1.48372);
\draw[pedge] (0.00000,0.00000) -- (12.00000,0.00000);
\draw[pedge] (4.04994,0.42777) -- (1.32113,1.43991);
\draw[pedge] (6.00000,10.39200) -- (1.32113,1.43991);
\draw[pedge] (6.00000,10.39200) -- (4.51516,5.61611);
\draw[pedge] (6.00000,10.39200) -- (8.02049,5.26608);
\draw[pedge] (12.00000,0.00000) -- (9.22826,0.75499);
\draw[pedge] (12.00000,0.00000) -- (10.01414,2.60858);
\draw[pedge] (2.69330,2.96621) -- (6.24482,1.48372);
\draw[pedge] (9.22826,0.75499) -- (6.24482,1.48372);
\draw[pedge] (4.51516,5.61611) -- (5.23909,3.48455);
\draw[pedge] (8.02049,5.26608) -- (7.28809,3.25717);
\draw[pedge] (6.18899,7.50100) -- (6.24482,1.48372);
\node[circle,fill=stateblue,inner sep=2.3pt] at (0.00000,0.00000) {};
\node[circle,fill=stateblue,inner sep=2.3pt] at (12.00000,0.00000) {};
\node[circle,fill=stateblue,inner sep=2.3pt] at (10.01414,2.60858) {};
\node[circle,fill=stateblue,inner sep=2.3pt] at (9.22826,0.75499) {};
\node[circle,fill=stateblue,inner sep=2.3pt] at (2.69330,2.96621) {};
\node[circle,fill=stateblue,inner sep=2.3pt] at (6.18899,7.50100) {};
\node[circle,fill=stateblue,inner sep=2.3pt] at (6.24482,1.48372) {};
\node[rectangle,fill=pairorange,inner sep=2.3pt] at (6.00000,10.39200) {};
\node[rectangle,fill=pairorange,inner sep=2.3pt] at (4.04994,0.42777) {};
\node[rectangle,fill=pairorange,inner sep=2.3pt] at (1.32113,1.43991) {};
\node[rectangle,fill=pairorange,inner sep=2.3pt] at (8.02049,5.26608) {};
\node[rectangle,fill=pairorange,inner sep=2.3pt] at (7.28809,3.25717) {};
\node[rectangle,fill=pairorange,inner sep=2.3pt] at (4.51516,5.61611) {};
\node[rectangle,fill=pairorange,inner sep=2.3pt] at (5.23909,3.48455) {};
\draw[line width=.7pt,fill=white] (0.00000,0.00000) circle[radius=2.4pt];
\draw[line width=.7pt,fill=white] (12.00000,0.00000) circle[radius=2.4pt];
\draw[line width=.7pt,fill=white] (6.00000,10.39200) circle[radius=2.4pt];
\node[anchor=north] at (6,-.25) {$t_0$ is the outer black face};
\end{tikzpicture}\\[1mm]Induced tree partition\\[4mm]\begin{tikzpicture}[x=.72cm,y=.72cm]
\path[use as bounding box] (-.6,-.7) rectangle (12.6,11);
\draw[base] (11.00000,5.00000) -- (8.98126,4.17880);
\draw[base] (11.00000,5.00000) -- (9.53553,1.46447);
\draw[base] (11.00000,5.00000) -- (9.53553,8.53553);
\draw[base] (8.98126,4.17880) -- (8.42742,2.55324);
\draw[base] (8.98126,4.17880) -- (7.51638,4.98317);
\draw[base] (9.53553,1.46447) -- (8.42742,2.55324);
\draw[base] (9.53553,1.46447) -- (6.00000,0.00000);
\draw[base] (9.53553,8.53553) -- (6.00000,10.00000);
\draw[base] (9.53553,8.53553) -- (7.52740,7.68875);
\draw[base] (8.42742,2.55324) -- (6.76545,2.01644);
\draw[base] (7.51638,4.98317) -- (7.06034,6.08414);
\draw[base] (7.51638,4.98317) -- (6.50752,4.68658);
\draw[base] (6.76545,2.01644) -- (6.00000,0.00000);
\draw[base] (6.76545,2.01644) -- (5.86894,3.49608);
\draw[base] (6.00000,0.00000) -- (2.46447,1.46447);
\draw[base] (6.00000,10.00000) -- (2.46447,8.53553);
\draw[base] (6.00000,10.00000) -- (5.98632,8.44657);
\draw[base] (7.52740,7.68875) -- (7.06034,6.08414);
\draw[base] (7.52740,7.68875) -- (5.98632,8.44657);
\draw[base] (2.46447,1.46447) -- (1.00000,5.00000);
\draw[base] (2.46447,1.46447) -- (3.16972,3.17873);
\draw[base] (1.00000,5.00000) -- (2.46447,8.53553);
\draw[base] (1.00000,5.00000) -- (2.71086,4.28652);
\draw[base] (3.16972,3.17873) -- (4.33384,3.78522);
\draw[base] (3.16972,3.17873) -- (2.71086,4.28652);
\draw[base] (2.46447,8.53553) -- (4.43156,7.65095);
\draw[base] (2.71086,4.28652) -- (3.96287,4.68083);
\draw[base] (4.43156,7.65095) -- (5.98632,8.44657);
\draw[base] (4.43156,7.65095) -- (4.84389,5.97076);
\draw[base] (7.06034,6.08414) -- (6.13725,5.58049);
\draw[base] (6.50752,4.68658) -- (5.86894,3.49608);
\draw[base] (6.50752,4.68658) -- (6.13725,5.58049);
\draw[base] (6.13725,5.58049) -- (4.84389,5.97076);
\draw[base] (5.86894,3.49608) -- (4.33384,3.78522);
\draw[base] (4.33384,3.78522) -- (3.96287,4.68083);
\draw[base] (4.84389,5.97076) -- (3.96287,4.68083);
\draw[ham] (11.00000,5.00000) -- (9.53553,1.46447);
\draw[ham] (11.00000,5.00000) -- (9.53553,8.53553);
\draw[ham] (8.98126,4.17880) -- (8.42742,2.55324);
\draw[ham] (8.98126,4.17880) -- (7.51638,4.98317);
\draw[ham] (8.42742,2.55324) -- (9.53553,1.46447);
\draw[ham] (9.53553,8.53553) -- (6.00000,10.00000);
\draw[ham] (6.00000,0.00000) -- (2.46447,1.46447);
\draw[ham] (6.00000,0.00000) -- (6.76545,2.01644);
\draw[ham] (2.46447,1.46447) -- (3.16972,3.17873);
\draw[ham] (1.00000,5.00000) -- (2.46447,8.53553);
\draw[ham] (1.00000,5.00000) -- (2.71086,4.28652);
\draw[ham] (2.46447,8.53553) -- (4.43156,7.65095);
\draw[ham] (6.00000,10.00000) -- (5.98632,8.44657);
\draw[ham] (7.51638,4.98317) -- (6.50752,4.68658);
\draw[ham] (7.52740,7.68875) -- (7.06034,6.08414);
\draw[ham] (7.52740,7.68875) -- (5.98632,8.44657);
\draw[ham] (7.06034,6.08414) -- (6.13725,5.58049);
\draw[ham] (6.76545,2.01644) -- (5.86894,3.49608);
\draw[ham] (6.50752,4.68658) -- (5.86894,3.49608);
\draw[ham] (6.13725,5.58049) -- (4.84389,5.97076);
\draw[ham] (3.16972,3.17873) -- (4.33384,3.78522);
\draw[ham] (4.33384,3.78522) -- (3.96287,4.68083);
\draw[ham] (4.43156,7.65095) -- (4.84389,5.97076);
\draw[ham] (2.71086,4.28652) -- (3.96287,4.68083);
\node[circle,draw,fill=black,inner sep=1.4pt] at (11.00000,5.00000) {};
\node[circle,draw,fill=white,inner sep=1.4pt] at (8.98126,4.17880) {};
\node[circle,draw,fill=white,inner sep=1.4pt] at (9.53553,1.46447) {};
\node[circle,draw,fill=white,inner sep=1.4pt] at (9.53553,8.53553) {};
\node[circle,draw,fill=black,inner sep=1.4pt] at (8.42742,2.55324) {};
\node[circle,draw,fill=black,inner sep=1.4pt] at (7.51638,4.98317) {};
\node[circle,draw,fill=white,inner sep=1.4pt] at (6.76545,2.01644) {};
\node[circle,draw,fill=black,inner sep=1.4pt] at (6.00000,0.00000) {};
\node[circle,draw,fill=black,inner sep=1.4pt] at (6.00000,10.00000) {};
\node[circle,draw,fill=black,inner sep=1.4pt] at (7.52740,7.68875) {};
\node[circle,draw,fill=white,inner sep=1.4pt] at (2.46447,1.46447) {};
\node[circle,draw,fill=black,inner sep=1.4pt] at (1.00000,5.00000) {};
\node[circle,draw,fill=black,inner sep=1.4pt] at (3.16972,3.17873) {};
\node[circle,draw,fill=white,inner sep=1.4pt] at (2.46447,8.53553) {};
\node[circle,draw,fill=white,inner sep=1.4pt] at (2.71086,4.28652) {};
\node[circle,draw,fill=black,inner sep=1.4pt] at (4.43156,7.65095) {};
\node[circle,draw,fill=white,inner sep=1.4pt] at (5.98632,8.44657) {};
\node[circle,draw,fill=white,inner sep=1.4pt] at (7.06034,6.08414) {};
\node[circle,draw,fill=white,inner sep=1.4pt] at (6.50752,4.68658) {};
\node[circle,draw,fill=black,inner sep=1.4pt] at (6.13725,5.58049) {};
\node[circle,draw,fill=black,inner sep=1.4pt] at (5.86894,3.49608) {};
\node[circle,draw,fill=white,inner sep=1.4pt] at (4.33384,3.78522) {};
\node[circle,draw,fill=white,inner sep=1.4pt] at (4.84389,5.97076) {};
\node[circle,draw,fill=black,inner sep=1.4pt] at (3.96287,4.68083) {};
\end{tikzpicture}\\[1mm]Hamiltonian cycle
\end{tabular}
\caption{The recovered partition includes the prescribed outer edge in one tree. The complementary dual cycle visits each of the 24 faces once and avoids the dual of that prescribed edge. This is the same cycle as in Figures 1 and 3.}
\end{figure}
\clearpage
\section*{13. gluing}
\begin{figure}[h!]\centering
\begin{tabular}{@{}c@{}}
\begin{tikzpicture}[x=.72cm,y=.72cm]
\draw[pedge] (0.00000,0.00000) -- (10.00000,0.00000);
\draw[base] (0.00000,0.00000) -- (5.00000,8.66000);
\draw[base] (0.00000,0.00000) -- (5.00000,1.73200);
\draw[pedge] (0.00000,0.00000) -- (4.00000,3.46400);
\draw[base] (10.00000,0.00000) -- (5.00000,8.66000);
\draw[base] (10.00000,0.00000) -- (5.00000,1.73200);
\draw[base] (10.00000,0.00000) -- (6.00000,3.46400);
\draw[base] (5.00000,8.66000) -- (4.00000,3.46400);
\draw[pedge] (5.00000,8.66000) -- (6.00000,3.46400);
\draw[base] (5.00000,1.73200) -- (4.00000,3.46400);
\draw[pedge] (5.00000,1.73200) -- (6.00000,3.46400);
\draw[base] (4.00000,3.46400) -- (6.00000,3.46400);
\node[circle,fill=stateblue,text=white,inner sep=2pt] at (0.00000,0.00000) {0};
\node[circle,fill=stateblue,text=white,inner sep=2pt] at (10.00000,0.00000) {1};
\node[rectangle,fill=pairorange,text=white,inner sep=2pt] at (5.00000,8.66000) {2};
\node[rectangle,fill=pairorange,text=white,inner sep=2pt] at (5.00000,1.73200) {3};
\node[circle,fill=stateblue,text=white,inner sep=2pt] at (4.00000,3.46400) {4};
\node[rectangle,fill=pairorange,text=white,inner sep=2pt] at (6.00000,3.46400) {5};
\end{tikzpicture}\\[1mm]First capped piece\\[4mm]\begin{tikzpicture}[x=.72cm,y=.72cm]
\draw[pedge] (0.00000,0.00000) -- (10.00000,0.00000);
\draw[base] (0.00000,0.00000) -- (5.00000,8.66000);
\draw[base] (0.00000,0.00000) -- (5.00000,1.73200);
\draw[pedge] (0.00000,0.00000) -- (4.00000,3.46400);
\draw[base] (10.00000,0.00000) -- (5.00000,8.66000);
\draw[base] (10.00000,0.00000) -- (5.00000,1.73200);
\draw[base] (10.00000,0.00000) -- (6.00000,3.46400);
\draw[base] (5.00000,8.66000) -- (4.00000,3.46400);
\draw[pedge] (5.00000,8.66000) -- (6.00000,3.46400);
\draw[base] (5.00000,1.73200) -- (4.00000,3.46400);
\draw[pedge] (5.00000,1.73200) -- (6.00000,3.46400);
\draw[base] (4.00000,3.46400) -- (6.00000,3.46400);
\node[circle,fill=stateblue,text=white,inner sep=2pt] at (0.00000,0.00000) {0};
\node[circle,fill=stateblue,text=white,inner sep=2pt] at (10.00000,0.00000) {1};
\node[rectangle,fill=pairorange,text=white,inner sep=2pt] at (5.00000,8.66000) {2};
\node[rectangle,fill=pairorange,text=white,inner sep=2pt] at (5.00000,1.73200) {6};
\node[circle,fill=stateblue,text=white,inner sep=2pt] at (4.00000,3.46400) {7};
\node[rectangle,fill=pairorange,text=white,inner sep=2pt] at (6.00000,3.46400) {8};
\end{tikzpicture}\\[1mm]Second capped piece
\end{tabular}
\caption{This separate nine-vertex witness glues two octahedra along vertices $0,1,2$. The blue trees intersect in edge $01$; the orange trees intersect in vertex $2$. Exhaustive checking verifies both induced trees in the glued graph. Identify the boundary vertices rather than placing both interiors in the same planar disk.}
\end{figure}
\clearpage
\section*{14. overview}
\begin{figure}[h!]\centering
\begin{tikzpicture}[x=1cm,y=1cm]
\tikzset{box/.style={draw,rounded corners=0pt,text width=10cm,align=center,minimum height=9mm}}
\node[box] (a) at (0,0) {Rooted triangulation $T$ and incidence matchings};
\node[box] (b) at (0,-1.5) {Disjoint pair $(r,s)$ and directed cycles $D$};
\node[box] (c) at (0,-3) {Cancellation + positive spoke circulation $\Rightarrow$ forest $Q_s$};
\node[box] (d) at (0,-4.5) {Retained spoke tree $\longleftrightarrow$ white-face cotree};
\node[box] (e) at (0,-6) {Selected edges $P$: one side of every face, no cycles};
\node[box] (f) at (0,-7.5) {Two induced trees $\longleftrightarrow$ Hamiltonian cycle in $G$};
\foreach \u/\v in {a/b,b/c,c/d,d/e,e/f} \draw[-{Stealth}] (\u)--(\v);
\end{tikzpicture}
\caption{Read down the main proof for a triangulation without separating triangles. The disk identity supports both the cancellation step and acyclicity of $P$. For a separating triangle, apply the construction to the capped pieces and glue using Figure 13. The flow records logical dependence, not an efficient algorithm.}
\end{figure}
\clearpage
\end{document}